\subsection{Pontryagin duality theorem}

THEOREM 2 (Pontryagin). --- The canonical mapping η from G into \(\widehat{G}\) is an isomorphism of topological groups. It maps the Haar measure dx to the bi-dual Haar measure d \(\widehat{x}\).

Let us first prove that \(\eta\) is injective and strict. It suffices for this to show that for every neighborhood U of \(e\) in G, there exists a neighborhood W of \(e\) in \(\widehat{\mathbf{G}}\) such that \(\overline{\eta}^{-1}(\mathrm{W})\subset \mathrm{U}\) (lemma 2 of II, p.~200). Now let V be a compact symmetric neighborhood of \(e\) in G such that \(\mathrm{V}^2\subset \mathrm{U}\), let \(f\) be a continuous positive function on G, with support contained in V, and such that \(f(e) > 0\). Let \(g = \widetilde{f}*f\). Then \(g\) belongs to \(\mathrm{A}(\mathrm{G})\), its support is contained in U, and \(g(e) > 0\). Moreover, \(\mathcal{F}_{\mathrm{G}}(g)\in \mathrm{L}^{1}(\widehat{\mathrm{G}})\) by prop. 11 of II, p.~217. The set W of \(\xi\) in \(\widehat{\mathbf{G}}\) such that

\[
\left| \overline {{\mathcal {F}}} _ {\widehat {\mathrm{G}}} (\mathcal {F} _ {\mathrm{G}} (g)) (\xi) - \overline {{\mathcal {F}}} _ {\widehat {\mathrm{G}}} (\mathcal {F} _ {\mathrm{G}} (g)) (e) \right| <   \frac {1}{2} g (e)
\]

is a neighborhood of \(e\) in \(\widehat{\mathbf{G}}\) since the function \(\overline{\mathcal{F}}_{\widehat{\mathrm{G}}}(\mathcal{F}_{\mathrm{G}}(g))\) is continuous on \(\widehat{\mathbf{G}}\) . Let \(x \in \overline{\eta}^{-1}(\mathrm{W})\) . By formula (26), we have

\[
\overline {{\mathcal {F}}} _ {\widehat {\mathrm{G}}} (\mathcal {F} _ {\mathrm{G}} (g)) (\eta (x)) = g (x),
\]

and therefore \(|g(x) - g(e)| < \frac{1}{2} g(e)\) . This implies \(g(x) \neq 0\) and therefore \(x \in \mathrm{U}\) , since the support of \(g\) is contained in U. Thus \(\overline{\eta}^{-1}(\mathrm{W}) \subset \mathrm{U}\) .

Let us prove that the map \(\eta\) is surjective. Since this map is a homeomorphism onto its image, the group \(\eta(\mathbf{G})\) is a locally compact subgroup of \(\widehat{\widehat{\mathbf{G}}}\). It is therefore closed in \(\widehat{\widehat{\mathbf{G}}}\) (TG, III, p.~22, cor. 2). Let us argue by contradiction and suppose that there exists a character \(\xi \in \widehat{\widehat{\mathbf{G}}}\) such that \(\xi \notin \eta(\mathbf{G})\). There then exists (corollary 1 of II, p.~219) a non-zero element \(f\) of \(\mathrm{L}^1(\widehat{\mathrm{G}})\) such that \(\mathcal{F}_{\widehat{\mathrm{G}}} (f)\) vanishes on \(\eta(\mathrm{G})\). Let \(g \in \mathrm{L}^1(\mathrm{G})\). The function \((x, \chi) \mapsto g(x)f(\chi)\overline{\langle\chi,x\rangle}\) belongs to \(\mathrm{L}^1(\mathrm{G} \times \widehat{\mathrm{G}})\). By the Lebesgue-Fubini theorem (INT, V, §8, no. 4, th. 1, a)), it follows that

\[
\begin{array}{r} \int_ {\widehat {\mathrm{G}}} f (\chi) \mathcal {F} _ {\mathrm{G}} (g) (\chi) d \chi = \int_ {\mathrm{G}} g (x) \Bigl (\int_ {\widehat {\mathrm{G}}} f (\chi) \overline {{\langle \chi , x \rangle}} d \chi \Bigr) d x \\ = \int_ {\mathrm{G}} g (x) \mathcal {F} _ {\widehat {\mathrm{G}}} (f) (\eta (x)) d x = 0. \end{array}
\]

Since the image of the Fourier transform is dense in \(\mathcal{C}_{0}(\widehat{\mathrm{G}})\) (cor. of prop. 5 of II, p.~209), it follows that the measure \(f \cdot d\chi\) is zero. This contradicts the fact that \(f \neq 0\) in \(\mathrm{L}^{1}(\widehat{\mathrm{G}})\), and proves that \(\eta\) is surjective.

The image measure \(\eta(dx)\) and the measure \(\nu\) dual to the measure \(d\chi\) are Haar measures on \(\widehat{\widehat{G}}\) . Let \(f\) be a non-zero element of A(G); in particular \(f \in \mathrm{L}^2(\mathrm{G})\) . By prop. 12 of II, p.~219, \(\mathcal{F}_{\mathrm{G}}(f) \in \mathrm{L}^{1}(\widehat{\mathrm{G}})\) and one has

\[
\int_ {\widehat {\widehat {\mathrm{G}}}} \Big | \overline {{\mathcal {F}}} _ {\widehat {\mathrm{G}}} (\mathcal {F} _ {\mathrm{G}} (f)) \Big | ^ {2} \eta (d x) = \int_ {\mathrm{G}} | f | ^ {2} d x = \int_ {\widehat {\widehat {\mathrm{G}}}} \Big | \overline {{\mathcal {F}}} _ {\widehat {\mathrm{G}}} (\mathcal {F} _ {\mathrm{G}} (f)) \Big | ^ {2} d \nu ,
\]

where the second equality follows from two applications of the Plancherel formula, so the dual Haar measure of \(d\chi\) is the measure \(\eta(dx)\) .

We shall henceforth identify G and \(\widehat{\mathrm{G}}\) by the isomorphism \(\eta\) . We then have:

COROLLARY. --- The Fourier cotransform of \(L^{2}(\widehat{G})\) on \(L^{2}(G)\) and the Fourier transform of \(L^{2}(G)\) on \(L^{2}(\widehat{G})\) are reciprocal isometries of each other.

Remark. --- Let \(f \in \mathrm{L}^2(\mathrm{G})\) and \(g \in \mathrm{L}^2(\widehat{\mathrm{G}})\) . Applying the Plancherel formula (24) to \(f\) and \(\overline{\mathcal{F}_{\widehat{\mathrm{G}}}(g)}\) , we obtain the formula

\[
\int_ {\mathrm{G}} f (x) \mathcal {F} _ {\widehat {\mathrm{G}}} (g) (x) d x = \int_ {\widehat {\mathrm{G}}} \mathcal {F} _ {\mathrm{G}} (f) (\chi) g (\chi) d \widehat {x} (\chi)\tag{29}
\]

since we have \(\mathcal{F}_{\mathrm{G}}(\overline{\mathcal{F}_{\widehat{\mathrm{G}}}(\boldsymbol{g})}) = \mathcal{F}_{\mathrm{G}}(\overline{\mathcal{F}}_{\widehat{\mathrm{G}}}(\overline{\boldsymbol{g}})) = \overline{\boldsymbol{g}}.\)

The Fourier transform and cotransform defined on \(\mathcal{M}^{1}(\widehat{\mathrm{G}})\) take values in the space of continuous bounded functions on G. For \(\beta\in\mathcal{M}^{1}(\widehat{\mathrm{G}})\) and \(x\in\mathrm{G}\) , one has

\[
\mathcal {F} _ {\widehat {\mathrm{G}}} (\beta) (x) = \int_ {\widehat {\mathrm{G}}} \overline {{\langle \chi , x \rangle}} d \beta (\chi), \qquad \overline {{\mathcal {F}}} _ {\widehat {\mathrm{G}}} (\beta) (x) = \int_ {\widehat {\mathrm{G}}} \langle \chi , x \rangle d \beta (\chi).
\]

The Fourier transforms of G and \(\widehat{G}\) are also transposes of each other. More precisely:

PROPOSITION 13. --- Let \(\alpha \in \mathcal{M}^{1}(\mathrm{G})\) and \(\beta \in \mathcal{M}^{1}(\widehat{\mathrm{G}})\) . We then have

\[
\mathcal {F} _ {\mathrm{G}} (\overline {{\mathcal {F}}} _ {\widehat {\mathrm{G}}} (\beta) \cdot \alpha) = \beta * \mathcal {F} _ {\mathrm{G}} (\alpha)\tag{30}
\]

and in particular

\[
\int_ {\mathrm{G}} \mathcal {F} _ {\widehat {\mathrm{G}}} (\beta) (x) d \alpha (x) = \int_ {\widehat {\mathrm{G}}} \mathcal {F} _ {\mathrm{G}} (\alpha) (\chi) d \beta (\chi).\tag{31}
\]

Formula (30) implies formula (31) by evaluating both sides of the identity at \(\chi = 1\) . Let us prove (30). Let \(\chi \in \widehat{G}\). It follows that

\[
\begin{array}{r l} & {(\mathcal {F} _ {\mathrm{G}} (\overline {{\mathcal {F}}} _ {\widehat {\mathrm{G}}} (\beta) \cdot \alpha)) (\chi) = \int_ {\mathrm{G}} \overline {{\langle \chi , x \rangle}} \overline {{\mathcal {F}}} _ {\widehat {\mathrm{G}}} (\beta) (x) d \alpha (x)} \\ & {\qquad = \int_ {\mathrm{G}} \overline {{\langle \chi , x \rangle}} \Big (\int_ {\widehat {\mathrm{G}}} \langle \xi , x \rangle d \beta (\xi) \Big) d \alpha (x).} \end{array}
\]

The function \((x,\xi)\mapsto\overline{\langle\chi,x\rangle}\langle\xi,x\rangle\) is continuous and bounded, hence integrable on \(G\times\widehat{G}\) with respect to the measure \(\alpha\otimes\beta\) . By the Lebesgue-Fubini theorem (INT, V, §8, no. 4, th. 1, a)), we obtain

\[
\begin{array}{r l} & {(\mathcal {F} _ {\mathrm{G}} (\overline {{\mathcal {F}}} _ {\widehat {\mathrm{G}}} (\beta) \cdot \alpha)) (\chi) = \int_ {\widehat {\mathrm{G}}} \Bigl (\int_ {\mathrm{G}} \overline {{\langle \chi \xi^ {- 1} , x \rangle}} d \alpha (x) \Bigr) d \beta (\xi)} \\ & {\qquad = \int_ {\widehat {\mathrm{G}}} \mathcal {F} _ {\mathrm{G}} (\alpha) (\chi \xi^ {- 1}) d \beta (\xi) = (\beta * \mathcal {F} _ {\mathrm{G}} (\alpha)) (\chi),} \end{array}
\]

as desired.

COROLLARY. --- The Fourier transform \(\mathcal{F}_{\mathrm{G}}\) is injective on \(\mathcal{M}^{1}(\mathrm{G})\) .

Indeed, if \(\alpha \in \mathcal{M}^1 (\mathrm{G})\) satisfies \(\mathcal{F}_{\mathrm{G}}(\alpha) = 0\), we deduce from (31) that \(\alpha (\mathcal{F}_{\widehat{\mathrm{G}}} (f)) = 0\) for every \(f\in \mathrm{L}^{1}(\widehat{\mathrm{G}})\). As the image of \(\mathrm{L}^1 (\widehat{\mathrm{G}})\) by the Fourier transform is dense in \(\mathcal{C}_0(\mathrm{G})\) (corollary of Prop. 5 of II, p.~209), we therefore have \(\alpha = 0\) .

There exist functional spaces on G and \(\widehat{\mathbf{G}}\) other than \(\mathrm{L}^2 (\mathrm{G})\) and \(\mathrm{L}^2 (\widehat{\mathrm{G}})\) on which \(\mathcal{F}\) and \(\overline{\mathcal{F}}\) are inverse isomorphisms of each other. The following theorem gives an example. We denote by B(G) the vector subspace of \(\mathrm{L}^1 (\mathrm{G})\) consisting of the elements \(f\in \mathrm{L}^1 (\mathrm{G})\) such that \(\mathcal{F}_{\mathrm{G}}(f)\in \mathrm{L}^1 (\widehat{\mathrm{G}})\). It is a subalgebra of \(\mathrm{L}^1 (\mathrm{G})\). Indeed, let \(f\) and \(g\) be in B(G). We have \(f*g\in \mathrm{L}^1 (\mathrm{G})\) and \(\mathcal{F}_{\mathrm{G}}(f*g) = \mathcal{F}_{\mathrm{G}}(f)\mathcal{F}_{\mathrm{G}}(g)\in \mathrm{L}^1 (\widehat{\mathrm{G}})\) since \(\mathcal{F}_{\mathrm{G}}(f)\in \mathrm{L}^1 (\widehat{\mathrm{G}})\) and \(\mathcal{F}_{\mathrm{G}}(g)\in \mathcal{C}_0(\widehat{\mathrm{G}})\) .

THEOREM 3. --- The restriction of the Fourier transform to B(G) induces an isomorphism of vector spaces from B(G) onto B( \(\widehat{G}\) ), whose inverse is induced by the restriction to B( \(\widehat{G}\) ) of the Fourier cotransform.

Let \(f \in \mathrm{B}(\mathrm{G})\) . Denote by \(g = \mathcal{F}_{\mathrm{G}}(f)\) . We have \(g \in \mathrm{L}^1(\widehat{\mathrm{G}}) \cap \mathcal{C}_0(\widehat{\mathrm{G}}) \subset \mathrm{L}^1(\widehat{\mathrm{G}}) \cap \mathrm{L}^2(\widehat{\mathrm{G}})\) . Set \(f_1 = \overline{\mathcal{F}}_{\widehat{\mathrm{G}}} (g) \in \mathrm{L}^2(\mathrm{G})\). For every continuous function with compact support \(h \in \mathcal{K}(\widehat{\mathrm{G}})\) , one has \(h \in \mathrm{L}^1(\widehat{\mathrm{G}}) \cap \mathrm{L}^2(\widehat{\mathrm{G}})\) and

\[
\begin{array}{r l} & {\int_ {\mathrm{G}} f _ {1} (x) \mathcal {F} _ {\widehat {\mathrm{G}}} (h) (x) d x = \int_ {\mathrm{G}} \overline {{\mathcal {F}}} _ {\widehat {\mathrm{G}}} (g) (x) \overline {{\overline {{\mathcal {F}}} _ {\widehat {\mathrm{G}}} (\overline {{h}}) (x)}} d x} \\ & {\qquad = \int_ {\widehat {\mathrm{G}}} g (\chi) h (\chi) d \widehat {x} (\chi)} \\ & {\qquad = \int_ {\widehat {\mathrm{G}}} \mathcal {F} _ {\mathrm{G}} (f) (\chi) h (\chi) d \widehat {x} (\chi) = \int_ {\mathrm{G}} f (x) \mathcal {F} _ {\widehat {\mathrm{G}}} (h) (x) d x} \end{array}
\]

using the Plancherel theorem and formula (31). Since \(\mathcal{K}(\widehat{\mathrm{G}})\) is dense in \(\mathrm{L}^2 (\widehat{\mathrm{G}})\) , its image under the Fourier transform is dense in \(\mathrm{L}^2 (\mathrm{G})\) . Consequently, we have \(f_{1} = f\) in \(\mathrm{L}^{1}(\mathrm{G})\) ; this shows that \(g\in \mathrm{B}(\widehat{\mathrm{G}})\) .

The formula \(f_{1}=f\) means that the restriction to B(G) of the composition \(\overline{F}_{\widehat{G}}\circ F_{G}\) is the identity map of B(G). Interchanging the roles of G and \(\widehat{G}\) , we observe that \(F_{G}\circ\overline{F}_{\widehat{G}}\) is the identity map of B( \(\widehat{G}\) ), which completes the proof of the theorem.

COROLLARY 1. --- Let \(f \in \mathrm{L}^{1}(\mathrm{G})\) . Then \(f \in \mathrm{B}(\mathrm{G})\) if and only if \(f\) belongs to the image of the Fourier transform \(\mathcal{F}_{\widehat{\mathrm{G}}}\) on \(\mathrm{L}^{1}(\widehat{\mathrm{G}})\) . In particular, we have \(\mathrm{A}(\mathrm{G}) \subset \mathrm{B}(\mathrm{G})\) .

Theorem 3 proves that if \(f \in \mathrm{B}(\mathrm{G})\) , then \(f = \mathcal{F}_{\widehat{\mathrm{G}}}(\overline{\mathcal{F}}_{\mathrm{G}}(f))\) , where \(\overline{\mathcal{F}}_{\mathrm{G}}(f)\) belongs to \(\mathrm{L}^1(\widehat{\mathrm{G}})\). Conversely, if \(f = \mathcal{F}_{\widehat{\mathrm{G}}}(g)\) , where \(g \in \mathrm{L}^1(\widehat{\mathrm{G}})\) , then we have \(g \in \mathrm{B}(\widehat{\mathrm{G}})\) and therefore \(f \in \mathrm{B}(\mathrm{G})\) by the theorem. The last assertion then follows from prop. 11 of II, p.~217.

COROLLARY 2. --- The vector space B(G) is an algebra both for multiplication and for convolution. The Fourier transform interchanges convolution and multiplication in B(G) and B( \(\widehat{G}\) ).

We have already seen that B(G) is a subalgebra of L \(^{1}\) (G). On the other hand, if \(f\) and \(g\) belong to B(G), then \(fg \in \text{L}^{1}(\text{G})\) since \(f \in \text{L}^{1}(\text{G})\) and \(g\) belongs to the image of the Fourier transform on L \(^{1}\) ( \(\widehat{\text{G}}\) ) (corollary 1). Since there exist \(f_{1}\) and \(g_{1}\) in L \(^{1}\) ( \(\widehat{\text{G}}\) ) such that \(f = \mathcal{F}_{\widehat{\text{G}}} (f_{1})\) and \(g = \mathcal{F}_{\widehat{\text{G}}} (g_{1})\) (loc. cit.), we have \(fg = \mathcal{F}_{\widehat{\text{G}}} (f_{1} * g_{1})\) , and hence \(fg \in \text{B(G)}\) again by the preceding corollary.

PROPOSITION 14. — Let \(f\) and \(g\) be in \(\mathrm{L}^2 (\mathrm{G})\) . Then \(\mathcal{F}_{\mathrm{G}}(fg) = \mathcal{F}_{\mathrm{G}}(f)*\mathcal{F}_{\mathrm{G}}(g)\) .

The equality holds if \(f\) and \(g\) belong to B(G) (Corollary 2), and in particular if \(f\) and \(g\) belong to A(G) since A(G) ⊂ B(G) (Cor. 1). Since A(G) is dense in L²(G) (Cor. 1 of II, p.~212), it suffices to show that the two members of the equality are continuous functions of \((f,g)\in\mathrm{L}^{2}(\mathrm{G})\times\mathrm{L}^{2}(\mathrm{G})\) with values in \(\mathcal{C}_{0}(\widehat{\mathrm{G}})\) . Now the map \((f,g)\mapsto\mathcal{F}_{\mathrm{G}}(fg)\) is obtained by composing the continuous map \((f,g)\mapsto fg\) of \(\mathrm{L}^{2}(\mathrm{G})\times\mathrm{L}^{2}(\mathrm{G})\) into \(\mathrm{L}^{1}(\mathrm{G})\) and the Fourier transform \(\mathcal{F}_{\mathrm{G}}\) of \(\mathrm{L}^{1}(\mathrm{G})\) into \(\mathcal{C}_{0}(\widehat{\mathrm{G}})\) , which is also continuous. Likewise, the map \((f,g)\mapsto\mathcal{F}_{\mathrm{G}}(f)*\mathcal{F}_{\mathrm{G}}(g)\) is obtained by composing the continuous maps \((f,g)\mapsto(\mathcal{F}_{\mathrm{G}}(f),\mathcal{F}_{\mathrm{G}}(g))\) of \(\mathrm{L}^{2}(\mathrm{G})\times\mathrm{L}^{2}(\mathrm{G})\) into \(\mathrm{L}^{2}(\widehat{\mathrm{G}})\times\mathrm{L}^{2}(\widehat{\mathrm{G}})\) and \((h_{1},h_{2})\mapsto h_{1}*h_{2}\) of \(\mathrm{L}^{2}(\widehat{\mathrm{G}})\times\mathrm{L}^{2}(\widehat{\mathrm{G}})\) into \(\mathcal{C}_{0}(\widehat{\mathrm{G}})\) (INT, VIII, §4, no. 5, prop. 15).

Remark. --- See no. 9 and exercises 22 of II, p.~270 and 31 of II, p.~275 for other examples of function spaces on which the Fourier transform is an isomorphism, in the case of particular groups G.

\subsection{Functorial properties of duality}

Let G and H be locally compact commutative groups. Recall that if \(\varphi\colon\mathbf{G}\to\mathbf{H}\) is a morphism of topological groups, the dual morphism \(\widehat{\varphi}\colon\widehat{\mathrm{H}}\to\widehat{\mathrm{G}}\) is defined by \(\langle\chi,\varphi(x)\rangle=\langle\widehat{\varphi}(\chi),x\rangle\) for all \(\chi\in\widehat{\mathrm{H}}\) and \(x\in\mathrm{G}\) . This definition shows that \(\widehat{\widehat{\varphi}}=\varphi\) with the identifications of G (resp. H) and \(\widehat{\widehat{\mathrm{G}}}\) (resp. \(\widehat{\widehat{\mathrm{H}}}\) ) of theorem 2 of II, p.~220.

Let \(\theta\) be a map from \(\mathbf{G} \times \mathbf{H}\) into \(\mathbf{U}\). For every \(x \in \mathbf{G}\) (respectively, every \(y \in \mathbf{H}\)), let \(\theta_x\) (respectively, \(\theta^y\)) be the function from \(\mathbf{G}\) into \(\mathbf{U}\) defined by \(y \mapsto \theta(x, y)\) (respectively, the function from \(\mathbf{H}\) into \(\mathbf{U}\) defined by \(x \mapsto \theta(x, y)\)). Suppose that the map \(\alpha: x \mapsto \theta_x\) is an isomorphism of the topological group \(\mathbf{G}\) onto the topological group \(\widehat{\mathbf{H}}\). For every \(y \in \mathbf{H}\) and \(x \in \mathbf{G}\), one has

\[
\theta^ {y} (x) = \theta (x, y) = \langle \alpha (x), y \rangle = \langle x, \widehat {\alpha} (y) \rangle ,
\]

that is, \(\theta^{y} = \widehat{\alpha}(y)\). By th. 2 of II, p.~220, the map \(\beta: y \mapsto \theta^{y}\) is therefore an isomorphism of the topological group H onto the topological group \(\widehat{G}\). Under these conditions, we shall say that \(\theta\) puts G and H in duality, or that G and H are in duality with respect to \(\theta\). We will then identify each of the groups G and H with the dual of the other. We will call the dual measure of the Haar measure dx the Haar measure on H obtained by transport of structure from the dual measure of dx.

Lemma 7. --- Let \((\mathrm{G}_{i})_{i\in\mathrm{I}}\) and \((\mathrm{H}_{i})_{i\in\mathrm{I}}\) be finite families of locally compact topological groups. For \(i\in I\) , let \(\theta_{i}\colon G_{i}\times H_{i}\to U\) be a map that puts the groups \(G_{i}\) and \(H_{i}\) in duality. The map \(\theta\) defined by

\[
\theta ((g _ {i}), (h _ {i})) = \prod_ {i \in I} \theta_ {i} (g _ {i}, h _ {i})
\]

puts the groups \(\prod G_{i}\) and \(\prod H_{i}\) in duality.

This follows from prop. 2 of II, p.~206 and from the definition preceding it.

DEFINITION 5. --- Let G, H, and K be topological groups. Let \(f: \mathrm{H} \to \mathrm{G}\) and \(g: \mathrm{G} \to \mathrm{K}\) be morphisms of topological groups. The pair \((f, g)\) is said to be an exact sequence of topological groups if it is an exact sequence of groups (A, II, p.~10, remark 5) and if \(f\) and \(g\) are strict morphisms.

We will represent an exact sequence by the diagram

\[
\mathrm{H} \xrightarrow {f} \mathrm{G} \xrightarrow {g} \mathrm{K},
\]

and we will say that a diagram

\[
\mathrm{G} _ {1} \xrightarrow {f _ {1}} \mathrm{G} _ {2} \xrightarrow {f _ {2}} \mathrm{G} _ {3} \rightarrow \dots \rightarrow \mathrm{G} _ {n} \xrightarrow {f _ {n}} \mathrm{G} _ {n + 1}
\]

is exact if each pair \((f_i, f_{i+1})\) for \(1 \leqslant i \leqslant n-1\) is exact.

A sequence

\[
1 \to \mathrm{H} \xrightarrow {f} \mathrm{G} \xrightarrow {g} \mathrm{K} \to 1
\]

is exact if and only if f is a strict injective morphism, g is a strict surjective morphism, and the kernel of g equals the image of f. If K is separated, the image of f is a closed subgroup of G.

Examples. --- 1) Let \(f\colon H \to G\) be a strict injective morphism whose image is a distinguished subgroup. The sequence

\[
1 \to \mathrm{H} \xrightarrow {f} \mathrm{G} \xrightarrow {p} \mathrm{G} / f (\mathrm{H}) \to 1,
\]

where p is the canonical projection, is exact. In particular, if H is a closed normal subgroup of G, the sequence of topological groups

\[
1 \to \mathrm{H} \xrightarrow {j} \mathrm{G} \xrightarrow {p} \mathrm{G} / \mathrm{H} \to 1,
\]

where j is the inclusion and p the canonical projection, is exact.

\begin{enumerate}
\def\labelenumi{\arabic{enumi})}
\setcounter{enumi}{1}
\tightlist
\item
  Let \(g: G \rightarrow K\) be a strict surjective morphism. The sequence
\end{enumerate}

\[
1 \to \mathrm{Ker} (g) \stackrel {j} {\longrightarrow} \mathrm{G} \stackrel {g} {\longrightarrow} \mathrm{K} \to 1,
\]

where j is the inclusion, is exact.

Theorem 4. --- A sequence

\[
\mathrm{H} \xrightarrow {f} \mathrm{G} \xrightarrow {g} \mathrm{K}
\]

of locally compact commutative topological groups is exact if and only if the dual sequence

\[
\widehat {\mathrm{K}} \xrightarrow {\widehat {g}} \widehat {\mathrm{G}} \xrightarrow {\widehat {f}} \widehat {\mathrm{H}}
\]

is exact.

We begin by proving a few lemmas. Note that each of them is, moreover, an easy consequence of the assertion of Thm. 4.

Lemma 8. --- Let \(g \colon G \to K\) be a morphism of locally compact commutative topological groups. If the morphism \(g\) is surjective and strict, then \(\widehat{g}\) is injective and strict.

Since \(g\) is surjective, the morphism \(\widehat{g}\) is injective (lemma 2 of II, p.~205). To prove that \(\widehat{g}\) is a strict morphism, it suffices to show that for every neighborhood U of \(e\) in \(\widehat{\mathbf{K}}\), there exists a neighborhood V of \(e\) in \(\widehat{\mathbf{G}}\) such that \(\overline{\widehat{g}}^{-1} (\mathrm{V})\subset \mathrm{U}\) (lemma 2 of II, p.~200). Let U be such a neighborhood of \(e\) in \(\widehat{\mathbf{K}}\). By definition of the topology of \(\widehat{\mathbf{K}}\), there exists a compact subset X of K and a number \(\varepsilon >0\) such that U contains the set of \(\widehat{z}\in \widehat{\mathbf{K}}\) which, for every \(z\in \mathrm{X}\), satisfy \(|\langle \widehat{z},z\rangle -1| < \varepsilon\). Since \(g\) is strict and surjective, there exists, by TG, I, p.~80, prop. 10, a compact subset \(\mathrm{X}_0\) of G such that \(g(\mathrm{X}_0) = \mathrm{X}\). Let V be the neighborhood of \(e\) in \(\widehat{\mathbf{G}}\) consisting of the elements \(\chi \in \widehat{\mathbf{G}}\) such that, for all \(x\in \mathrm{X}_0\), one has \(|\langle \chi ,x\rangle -1| < \varepsilon\). We then have \(\overline{\widehat{g}}^{-1} (\mathrm{V})\subset \mathrm{U}\). This proves the assertion.

Lemma 9. --- Let \(f\colon H\to G\) be a morphism of locally compact topological groups. If the morphism \(f\) is injective and strict, then \(\widehat{f}\) is surjective and strict.

Suppose that \(f\) is injective and strict. The morphism \(\widehat{f}\) induces, by passing to the quotient, a morphism \(q \colon \widehat{\mathrm{G}} / \mathrm{Ker}(\widehat{f}) \to \widehat{\mathrm{H}}\). It remains to show that this is an isomorphism of topological groups; by duality, it suffices for this to show that its dual \(\widehat{q}\) is an isomorphism.

Let L denote the dual group of \(\widehat{\mathrm{G}}/\mathrm{Ker}(\widehat{f})\) and \(p\colon\widehat{\mathrm{G}}\to\widehat{\mathrm{G}}/\mathrm{Ker}(\widehat{f})\) the canonical projection. We shall first prove that \(\widehat{p}\) induces, by passing to subspaces, an isomorphism from L onto \(f(\mathrm{H})\) .

We have \(q \circ p = \widehat{f}\) , whence \(\widehat{p} \circ \widehat{q} = f\) . The image of \(\widehat{p}\) therefore contains \(f(\mathrm{H})\) .

Since \(f\) is strict, its image \(f(\mathrm{H})\) is a locally compact subgroup of \(\mathrm{G}\) , and is therefore closed (TG, III, p.~22, cor. 2). Let \(\mathrm{K} = \mathrm{G}/f(\mathrm{H})\) and consider the exact sequence

\[
1 \to \mathrm{H} \xrightarrow {f} \mathrm{G} \xrightarrow {g} \mathrm{K} \to 1
\]

associated with example 1. By duality, the morphism \(\widehat{f} \circ \widehat{g}\) is trivial, and hence the image of \(\widehat{g}\) is contained in \(\operatorname{Ker}(\widehat{f}) = \operatorname{Ker}(p)\) . Thus \(p \circ \widehat{g}\) is the trivial morphism and, again by duality, \(g \circ \widehat{p}\) is also trivial. It follows that the image of \(\widehat{p}\) is contained in the kernel of \(g\) , which is equal to \(f(\mathrm{H})\) . We conclude that the image of \(\widehat{p}\) is equal to \(f(\mathrm{H})\) .

Moreover, since p is a surjective and strict morphism, the dual morphism \(\widehat{p}\) is a strict injective morphism from L to G (Lemma 8). It follows that \(\widehat{p}\) induces an isomorphism of topological groups from L onto \(f(\mathrm{H})\). Since \(\widehat{p} \circ \widehat{q} = f\) and since f induces an isomorphism from H onto its image \(f(\mathrm{H})\), the morphism \(\widehat{q}\) is an isomorphism.

Lemma 10. --- Let

\[
\mathrm{H} \xrightarrow {f} \mathrm{G} \xrightarrow {g} \mathrm{K}
\]

an exact sequence of locally compact abelian groups. The kernel of \(\hat{f}\) is equal to the image of \(\hat{g}\) .

The homomorphism \(\widehat{f} \circ \widehat{g}\) is trivial by duality, so the image of \(\widehat{g}\) is contained in the kernel of \(\widehat{f}\). Conversely, let \(\chi\) be in the kernel of \(\widehat{f}\). This means that \(\operatorname{Im}(f) = \operatorname{Ker}(g)\) is contained in the kernel of \(\chi\), hence that there exists a character \(\eta\) of \(\operatorname{Im}(g)\) such that \(\eta \circ g = \chi\). Since the inclusion of \(\operatorname{Im}(g)\) in K is strict, the dual map of restriction of characters from K to \(\operatorname{Im}(g)\) is surjective (Lemma 9). There therefore exists a character \(\beta\) of K such that \(\eta\) is the restriction of \(\beta\), and it follows that \(\chi = \beta \circ g = \widehat{g}(\beta)\). From this we conclude that the kernel of \(\widehat{f}\) is contained in the image of \(\widehat{g}\).

Now let us prove Theorem 4. By duality, it suffices to show that the sequence \(\widehat{\mathrm{K}}\xrightarrow{\widehat{g}}\widehat{\mathrm{G}}\xrightarrow{\widehat{f}}\widehat{\mathrm{H}}\) is exact when the sequence \(\mathrm{H}\xrightarrow{f}\mathrm{G}\xrightarrow{g}\mathrm{K}\) is exact. Now, by Lemmas 8 and 9, the morphisms \(\widehat{f}\) and \(\widehat{g}\) are strict and, by Lemma 10, the kernel of \(\widehat{f}\) is equal to the image of \(\widehat{g}\).

COROLLARY 1. --- Let

\[
1 \to \mathrm{H} \xrightarrow {f} \mathrm{G} \xrightarrow {g} \mathrm{K} \to 1
\]

an exact sequence of locally compact commutative topological groups. The morphism \(\widehat{g}\) induces an isomorphism between \(\widehat{\mathrm{K}}\) and \(f(\mathrm{H})^{\perp}\) , and \(\widehat{f}\) induces, by passing to the quotient, an isomorphism between \(\widehat{\mathrm{G}} / f(\mathrm{H})^{\perp}\) and \(\widehat{\mathrm{H}}\) .

By Theorem 4, the sequence

\[
1 \to \widehat {\mathrm{K}} \xrightarrow {\widehat {g}} \widehat {\mathrm{G}} \xrightarrow {\widehat {f}} \widehat {\mathrm{H}} \to 1\tag{32}
\]

is exact. The morphism \(\widehat{g}\) therefore induces an isomorphism from \(\widehat{\mathbf{K}}\) onto \(\mathrm{Ker}(\widehat{f}) = f(\mathrm{H})^{\perp}\) (lemma 2 of II, p.~205), and \(\widehat{f}\) induces, by passing to the quotient, an isomorphism from \(\widehat{\mathrm{G}} / \mathrm{Ker}(\widehat{f}) = \widehat{\mathrm{G}} / f(\mathrm{H})^{\perp}\) onto \(\widehat{\mathrm{H}}\) (loc. cit.).

COROLLARY 2. --- Let \(f\colon G\to H\) be a morphism of locally compact commutative topological groups. The morphism \(f\) is strict if and only if \(\widehat{f}\) is strict.

By the canonical decomposition (E, II, p.~44) of a strict morphism, this follows from Lemmas 8 and 9.

COROLLARY 3. --- Let H be a subgroup of G. Then (H \(^{\perp}\) ) \(^{\perp}\) = \(\overline{H}\) .

Since \(H^{\perp} = \overline{H}^{\perp}\) , we may assume that H is closed. Let \(f: H \to G\) be the canonical injection and \(p: G \to G/H\) the canonical projection. We have \(H^{\perp} = \operatorname{Ker}(\widehat{f})\) (Lemma 10). Let k be the canonical injection of \(H^{\perp}\) into \(\widehat{G}\). By Theorem 4, the morphism \(\widehat{p}\) induces an isomorphism \(\widehat{p}_{H}: \widehat{G/H} \to H^{\perp}\) of topological groups and we have \(k \circ \widehat{p}_{H} = \widehat{p}\). Consequently (Corollary 1), we obtain

\[
(\mathrm{H} ^ {\perp}) ^ {\perp} = \mathrm{Ker} (\widehat {k}) = \mathrm{Ker} (\widehat {\widehat {p}} _ {\mathrm{H}} \circ \widehat {k}) = \mathrm{Ker} (p) = \mathrm{H}.
\]

COROLLARY 4. --- Let I be a set and let \((\mathrm{H}_{i})_{i\in\mathrm{I}}\) be a family of closed subgroups of G. The orthogonal of the closed subgroup generated by the \(H_{i}\) is \(\bigcap_{i\in I}H_{i}^{\perp}\) . The orthogonal of \(\bigcap_{i}H_{i}\) is the closed subgroup generated by the subgroups \(H_{i}^{\perp}\) .

The first assertion follows from the definition of the orthogonal. Applying this result and Cor. 3 to the family of closed subgroups \((\mathrm{H}_{i}^{\perp})_{i\in\mathrm{I}}\) of \(\widehat{G}\) we see that \(\bigcap_{i}H_{i}\) is the orthogonal of the closed subgroup generated by the subgroups \(H_{i}^{\perp}\) and the second assertion is then obtained by duality.

COROLLARY 5. --- Let \(\varphi\) : G → H be a morphism of locally compact commutative groups. Then the subgroup \(\overline{\text{Im}(\varphi)}\) of H and the subgroup \(\text{Ker}(\widehat{\varphi})\) of \(\widehat{\text{H}}\) are orthogonal to each other. In particular, for \(\widehat{\varphi}\) to be injective, it is necessary and sufficient that the image of \(\varphi\) be dense in H.

We have \(\operatorname{Ker}(\widehat{\varphi}) = \varphi(\mathrm{G})^{\perp}\) (lemma 2 of II, p.~205), whence the result by cor. 3.

COROLLARY 6. --- Let \(k \in \mathbf{Z}\). Then the kernel of the homomorphism \(x \mapsto x^k\) from G to G and the closure of the image of the morphism \(\chi \mapsto \chi^k\) from \(\widehat{\mathrm{G}}\) to \(\widehat{\mathrm{G}}\) are orthogonal to each other.

This follows from the preceding corollary, since the morphisms \(x \mapsto x^{k}\) from G to G and \(\chi \mapsto \chi^{k}\) from \(\widehat{G}\) to \(\widehat{G}\) are dual to each other.

Recall (A, X, p.~17) that a commutative group A is divisible if, for every nonzero \(n \in \mathbf{Z}\), the map \(a \mapsto a^n\) from A to A is surjective.

COROLLARY 7. --- Let G be a locally compact abelian group.

\begin{enumerate}
\def\labelenumi{\alph{enumi})}
\item
  If G is divisible, then the dual group \(\widehat{\mathrm{G}}\) is torsion-free;
\item
  If the dual group \(\widehat{\mathbf{G}}\) is torsion-free, and if \(k\in\mathbf{Z}\) is nonzero, then the image of the homomorphism \(x\mapsto x^{k}\) from G to G is dense in G;
\item
  Suppose G is discrete or compact. For G to be divisible it is necessary and sufficient that \(\widehat{G}\) be torsion-free.
\end{enumerate}

The assertions \(a\) ) and \(b\) ) follow from cor. 6. If G is discrete or compact, the image of the homomorphism \(x \mapsto x^k\) from G to G is closed, and \(c\) ) follows from \(a\) ) and \(b\) ).

Remark. --- There exist locally compact abelian groups G that are not divisible and such that \(\widehat{G}\) is torsion-free (exercise 63 of II, p.~299).

\subsection{The Poisson formula}

In this section, we consider a closed subgroup H of G. We denote \(\beta = dx\) the Haar measure on G and \(\widehat{\beta}\) the dual Haar measure on \(\widehat{G}\) . One denotes by \(\alpha\) a Haar measure on H and \(\widehat{\alpha}\) the dual Haar measure on the dual group \(\widehat{H}\) , which we identify with \(\widehat{G}/H^{\perp}\) (Theorem 4 of

II, p.~226). We also identify \(\widehat{\mathrm{G}}/\widehat{\mathrm{H}}\) with \(\mathrm{H}^{\perp}\) via the dual mapping of the canonical projection \(\mathrm{G}\to\mathrm{G}/\mathrm{H}\) (loc. cit.).

We shall denote by \(\dot{x}\) the canonical image of an element \(x\) of G in G/H and by \(\dot{\chi}\) the canonical image of an element \(\chi\) of \(\widehat{G}\) in \(\widehat{G}/H^{\perp}\) .

We denote by \(\gamma\) the Haar measure \(\beta/\alpha\) on G/H (INT, VII, §2, no. 2, def. 1 and no. 7, prop. 10), and \(\widehat{\gamma}\) the dual Haar measure on \(H^{\perp}\). Recall (INT, VII, §2, n° 3, prop. 5, c)) that the measure \(\gamma\) is characterized by the following property: for every \(f \in \mathcal{L}^{1}(G)\), the function \(y \mapsto f(xy)\) on H is \(\alpha\)-integrable for \(\beta\)-almost all \(x \in G\); its integral depends only on \(\dot{x}\) and the function defined \(\gamma\)-almost everywhere on G/H by

\[
f ^ {\flat} \colon \dot {x} \mapsto \int_ {\mathrm{H}} f (x h) d \alpha (h)
\]

belongs to L \(^{1}\) (G/H, γ) and satisfies

\[
\int_ {\mathrm{G} / \mathrm{H}} f ^ {\flat} d \gamma = \int_ {\mathrm{G}} f d \beta .\tag{33}
\]

PROPOSITION 15. --- Let \(f \in \mathrm{L}^1(\mathrm{G})\) such that the restriction to \(\mathrm{H}^\perp\) of the continuous function \(\mathcal{F}_{\mathrm{G}}(f)\) is integrable with respect to \(\widehat{\gamma}\). Then, for almost all \(x \in \mathrm{G}\), the function \(y \mapsto f(xy)\) on \(\mathrm{H}\) is \(\alpha\)-integrable, and one has:

\[
\int_ {\mathrm{H}} f (x y) d \alpha (y) = \int_ {\mathrm{H} ^ {\perp}} \langle \chi , x \rangle \mathcal {F} _ {\mathrm{G}} (f) (\chi) d \widehat {\gamma} (\chi).
\]

Based on the foregoing, the function \(f^{\flat}\) defined almost everywhere on G/H by

\[
f ^ {\flat} (\dot {x}) = \int_ {\mathrm{H}} f (x y) d \alpha (y)
\]

belongs to \(L^{1}(G/H)\) . The Fourier transform of \(f^{\flat}\) is identified with the function on \(H^{\perp} = \widehat{G/H}\) given for \(\chi \in H^{\perp}\) by

\[
\begin{array}{r} \mathcal {F} _ {\mathrm{G/H}} (f ^ {\flat}) (\chi) = \int_ {\mathrm{G/H}} \overline {{\langle \chi , \dot {x} \rangle}} f ^ {\flat} (\dot {x}) d \gamma (\dot {x}) \\ = \int_ {\mathrm{G}} \overline {{\langle \chi , x \rangle}} f (x) d \beta (x) = \mathcal {F} _ {\mathrm{G}} (f) (\chi) \end{array}
\]

by formula (33), applied to the integrable function \(x \mapsto \overline{\langle\chi, x\rangle} f(x)\) . By hypothesis, the function \(\mathcal{F}(f)|\mathrm{H}^{\perp} = \mathcal{F}_{\mathrm{G}/\mathrm{H}}(f^{\flat})\) belongs to \(\mathrm{L}^{1}(\mathrm{H}^{\perp})\) , and hence the function \(f^{\flat}\) belongs to the space B(G/H). It follows (th. 3 of II, p.~222) that \(f^{\flat}\) coincides almost everywhere with \(\overline{\mathcal{F}}_{\widehat{\mathrm{G / H}}}(\mathcal{F}_{\mathrm{G / H}}(f^{\flat}))\) . For almost every \(\dot{x}\in \mathrm{G / H}\) , we therefore have

\[
f ^ {\flat} (\dot {x}) = \int_ {\mathrm{H} ^ {\perp}} \langle \chi , x \rangle \mathcal {F} _ {\mathrm{G/H}} (f ^ {\flat}) (\chi) d \widehat {\gamma} (\chi) = \int_ {\mathrm{H} ^ {\perp}} \langle \chi , x \rangle \mathcal {F} _ {\mathrm{G}} (f) (\chi) d \widehat {\gamma} (\chi).
\]

This concludes the proof.

COROLLARY (Poisson formula). --- Let \(f \in \mathcal{L}^{1}(G)\) . It is assumed that the following conditions are satisfied:

\begin{enumerate}
\def\labelenumi{(\roman{enumi})}
\item
  The restriction of \(\mathcal{F}_{\mathrm{G}}(f)\) to \(\mathrm{H}^{\perp}\) is integrable;
\item
  For every \(x \in G\) , the function \(y \mapsto f(xy)\) on H is integrable;
\item
  The map \(x\mapsto \int_{\mathrm{H}}f(xy)d\alpha (y)\) is continuous on G. Then one has
\end{enumerate}

\[
\int_ {\mathrm{H}} f (y) d \alpha (y) = \int_ {\mathrm{H} ^ {\perp}} \mathcal {F} _ {\mathrm{G}} (f) (\chi) d \widehat {\gamma} (\chi).\tag{34}
\]

Indeed, taking up the notation of the proof of the preceding proposition, the functions \(f^{\flat}\) and \(\overline{\mathcal{F}}_{\widehat{\mathrm{G/H}}}\left(\mathcal{F}_{\mathrm{G/H}}(f^{\flat})\right)\) on G/H are continuous and equal almost everywhere. They are therefore equal everywhere and in particular at e, which gives formula (34).

PROPOSITION 16. --- The measure \(\widehat{\alpha}\) on \(\widehat{\mathrm{H}} = \widehat{\mathrm{G}} / \mathrm{H}^{\perp}\) is equal to \(\widehat{\beta} / \widehat{\gamma}\) .

Fix \(f\in \mathcal{K}(\mathrm{G})\) nonzero. For \(x\in \mathrm{G}\) and \(\chi \in \widehat{\mathrm{G}}\) , set

\[
\varphi (x, \chi) = \int_ {\mathrm{H}} f (x y) \langle \chi , y \rangle d \alpha (y).
\]

The function \(\varphi\) is continuous on \(G \times \widehat{G}\) (INT, IV, §4, no. 3, cor. 1 of th. 2). For \(x\) fixed, \(\varphi(x, \chi)\) depends only on the class of \(\chi\) in \(\widehat{G}/H^{\perp} = \widehat{H}\) . For \(\chi\) fixed, \(\langle \chi, x \rangle \varphi(x, \chi)\) depends only on the class of \(x\) in \(G/H\) , and the function \(\dot{x} \mapsto \langle \chi, x \rangle \varphi(x, \chi)\) on \(G/H\) has compact support.

Let \(x \in G\) . The function \(\dot{\chi} \mapsto \varphi(x, \chi)\) on \(\hat{H}\) is the Fourier cotransform of the function \(y \mapsto f(xy)\) on H. The latter is square-integrable, so by the Plancherel formula (23) of II, p.~217, we have

\[
\int_ {\widehat {\mathrm{G}} / \mathrm{H} ^ {\perp}} | \varphi (x, \chi) | ^ {2} d \widehat {\alpha} (\dot {\chi}) = \int_ {\mathrm{H}} | f (x y) | ^ {2} d \alpha (y).\tag{35}
\]

 Let \(\chi\in\widehat{\mathrm{G}}\) . The function \(\dot{x}\mapsto\langle\chi,x\rangle\varphi(x,\chi)\) belongs to \(\mathcal{H}(\mathrm{G}/\mathrm{H})\) , hence to \(\mathrm{L}^{1}(\mathrm{G}/\mathrm{H})\) . Its Fourier cotransform is the function on \(\mathrm{H}^{\perp}\)

whose value at \(\xi\in\mathrm{H}^{\perp}\) is

\[
\begin{array}{r l} & {\int_ {\mathrm{G/H}} \langle \xi , \dot {x} \rangle \langle \chi , x \rangle \varphi (x, \chi) d \gamma (\dot {x}) = \int_ {\mathrm{G/H}} \Bigl (\int_ {\mathrm{H}} \langle \chi \xi , x y \rangle f (x y) d \alpha (y) \Bigr) d \gamma (\dot {x})} \\ & {\qquad = \int_ {\mathrm{G}} \langle \chi \xi , x \rangle f (x) d \beta (x) = \overline {{\mathcal {F}}} _ {\mathrm{G}} (f) (\chi \xi)} \end{array}
\]

by formula (33). Therefore

\[
\int_ {\mathrm{G/H}} | \varphi (x, \chi) | ^ {2} d \gamma (\dot {x}) = \int_ {\mathrm{H} ^ {\perp}} | \overline {{\mathcal {F}}} _ {\mathrm{G}} (f) (\chi \xi) | ^ {2} d \widehat {\gamma} (\xi)\tag{36}
\]

by the Plancherel formula again.

One then finally computes

\[
\begin{array}{l l} \int_ {\widehat {G}} | \overline {{\mathcal {F}}} _ {G} (f) | ^ {2} d \widehat {\beta} = \int_ {G} | f | ^ {2} d \beta & (\text {by (23)}) \\ = \int_ {G / H} d \gamma (\dot {x}) \int_ {H} | f (x y) | ^ {2} d \alpha (y) & (\text {by (33)}) \\ = \int_ {G / H} d \gamma (\dot {x}) \int_ {\widehat {G} / H ^ {\perp}} | \varphi (x, \chi) | ^ {2} d \widehat {\alpha} (\dot {\chi}) & (\text {by (35)}) \\ = \int_ {\widehat {G} / H ^ {\perp}} d \widehat {\alpha} (\dot {\chi}) \int_ {G / H} | \varphi (x, \chi) | ^ {2} d \gamma (\dot {x}) \\ = \int_ {\widehat {G} / H ^ {\perp}} d \widehat {\alpha} (\dot {\chi}) \int_ {H ^ {\perp}} | \overline {{\mathcal {F}}} _ {G} (f) (\chi \xi) | ^ {2} d \widehat {\gamma} (\xi) & (\text {by (36)}) \end{array}
\]

where INT, V, §8, no. 3, prop. 5 has been applied to the continuous positive function \((\dot{x},\dot{\chi})\mapsto |\varphi (x,\chi)|^{2}\) on \(\mathrm{G / H}\times \widehat{\mathrm{G}} /\mathrm{H}^{\perp}\) .

Comparing this equality with the integration formula (33) for the group \(\widehat{G}\) , one concludes that the Haar measures \(\widehat{\alpha}\) and \(\widehat{\beta}/\widehat{\gamma}\) coincide.

\subsection{Examples of duality}

PROPOSITION 17. --- Let \(n \geqslant 1\) be an integer. Denote by \(\boldsymbol{\mu}_{n}\) the group of \(n\) -th roots of unity in \(\mathbf{C}\) . The groups \(\mathbf{Z}/n\mathbf{Z}\) and \(\boldsymbol{\mu}_{n}\) are in duality with respect to the map induced by passage to the quotient of the map \(\mathbf{Z} \times \boldsymbol{\mu}_{n} \to \mathbf{U}\) defined by \((m,z) \mapsto z^{m}\) .

The group \(\mathbf{Z}/n\mathbf{Z}\) coincides with the set of homomorphisms \(\chi\) from \(\mathbf{Z}/n\mathbf{Z}\) to \(\mathbf{U}\) . These are of the form \(m \mapsto \chi(1)^m\) where \(\chi(1)\) is an arbitrary element of \(\mathbf{U}\) such that \(\chi(1)^n = 1\) , whence the result.

COROLLARY 1. --- Let G be a finite commutative group. The dual group \(\widehat{\mathbf{G}}\) is isomorphic to G.

The group G is isomorphic to a finite product of cyclic groups (A, VII, p.~22, th. 3), and its dual group is isomorphic to the product of the dual groups of these (prop. 2 of II, p.~206). One is therefore reduced to the case where G is cyclic, which follows from proposition 17 since the group \(\mu_{n}\) is cyclic of order n (A, V, p.~75, th. 1).

COROLLARY 2. --- Let G be a locally compact commutative group. The group G is finite if and only if \(\widehat{\mathbf{G}}\) is finite. A closed subgroup H of G has finite index if and only if its orthogonal is finite.

By duality, the first assertion follows from the fact that the dual of a finite group is finite (corollary 1). The second follows from it, since \(\widehat{\mathrm{G / H}}\) is identified with \(\mathrm{H}^{\perp}\) (th. 4 of II, p.~226).

PROPOSITION 18. --- For \(\widehat{\mathbf{G}}\) to be compact, it is necessary and sufficient that \(\mathbf{G}\) be discrete. If \(\mathbf{G}\) is discrete, the dual measure of the counting measure on \(\mathbf{G}\) is the normalized Haar measure on \(\widehat{\mathbf{G}}\) . If \(\mathbf{G}\) is compact, the dual measure of the normalized Haar measure is the counting measure on \(\widehat{\mathbf{G}}\) .

Suppose G is discrete, and let \(\alpha\) be the counting measure on G. Let \(\varphi\) be the characteristic function of \(e \in G\) . We have \(\mathcal{F}_{\mathrm{G}}(\varphi) = 1\) on \(\widehat{G}\) . Since \(\mathcal{F}_{\mathrm{G}}(\varphi)\) tends to 0 at infinity, the group \(\widehat{G}\) is compact.

Furthermore, for the dual measure \(\widehat{\alpha}\) of \(\alpha\), the function \(\mathcal{F}_{\mathrm{G}}(\varphi)\) must have integral \(\varphi(e)=1\) (prop. 12 of II, p.~219). Therefore \(\widehat{\alpha}(\widehat{\mathrm{G}})=1\) .

Suppose G is compact. Then the Haar measure dx belongs to \(\mathcal{M}^{1}(G)\) . Its Fourier transform is strictly positive at \(\chi = e\) and zero for \(\chi \neq 0\) (prop. 6 of II, p.~210). Since it is continuous on \(\widehat{G}\) , the group \(\widehat{G}\) is discrete. If the measure of G is 1, we deduce by duality from the previous case that the dual measure of the measure dx is the counting measure on \(\widehat{G}\) .

COROLLARY 1 (Orthogonality relations). --- Suppose G is discrete and equipped with the counting measure. For x and y in G, we have

\[
\int_ {\widehat {G}} \chi (x) \overline {{\chi (y)}} d \chi = \left\{ \begin{array}{l l} 0 & \text{if } x \neq y \\ 1 & \text{if } x = y. \end{array} \right.
\]

This follows from the cor. of theo. 1 of II, p.~215 and from duality.

COROLLARY 2. --- Let H be a closed subgroup of G.

\begin{enumerate}
\def\labelenumi{\alph{enumi})}
\item
  For H to be compact, it is necessary and sufficient that \(H^{\perp}\) be open in \(\widehat{G}\) ;
\item
  For H to be open, it is necessary and sufficient that \(H^{\perp}\) be compact in \(\widehat{G}\) .
\item
  To say that \(H^{\perp}\) is open amounts to saying that \(\widehat{G}/H^{\perp}\) is discrete, but \(\widehat{G}/H^{\perp}\) is isomorphic to \(\widehat{H}\) (th. 4 of II, p.~226); the assertion therefore follows from prop. 18. Assertion b) follows by duality from assertion a) applied to \(H^{\perp}\) .
\end{enumerate}

COROLLARY 3. --- Let \((\mathrm{H}_{i})_{i\in\mathrm{I}}\) be a decreasing filtered family of compact subgroups of G. For G to be identified with the projective limit of the groups \(G/H_{i}\) , it is necessary and sufficient that \(\hat{G}\) be the union of the open subgroups \(H_{i}^{\perp}\) .

To say that G is identified with the projective limit of the \(G/H_{i}\) amounts to saying that \(\bigcap_{i}H_{i}=\{e\}\) (TG, III, p.~60, prop. 2), that is, that \(\bigcup_{i}H_{i}^{\perp}\) is dense in \(\widehat{G}\) (cor. 4 of th. 4 of II, p.~226). Now \(\bigcup_{i}H_{i}^{\perp}\) is an open, hence closed, subgroup of \(\widehat{G}\) .

COROLLARY 4. --- Let I be a set and let \((\mathrm{H}_{i})_{i\in\mathrm{I}}\) be a family of compact groups. The dual of the product group of \(H_{i}\) is the discrete group direct sum of the groups \(\widehat{H}_{i}\) .

This is a special case of Cor. 3.

PROPOSITION 19. --- Let K be a locally compact, non-discrete field, not necessarily commutative, and let G be the additive group of K, whose group law is written additively. Let \(\chi\) be a unitary character of G distinct from 1. For \(x,y\in\mathbf{G}\), set \(\theta(x,y)=\chi(xy)\). Then G is in duality with itself with respect to \(\theta\) .

For \(y \in \mathbf{G}\), let \(\chi_y\) be the map from \(\mathbf{G}\) into \(\mathbf{U}\) such that \(\chi_y(x) = \chi(xy)\). We have \(\chi_y \in \widehat{\mathbf{G}}\), and it must be shown that \(\beta: y \mapsto \chi_y\) is a topological group isomorphism from \(\mathbf{G}\) into \(\widehat{\mathbf{G}}\).

The map \(\beta\) is an injective homomorphism from G into \(\widehat{\mathbf{G}}\); it is continuous (TG, X, p.~28, th. 3 applied to the continuous map \(\theta\) of \(\mathbf{G} \times \mathbf{G}\) into \(\mathbf{C}\)). Let us prove that \(\theta\) is a homeomorphism onto its image. It suffices (Lemma 2 of II, p.~200) to show that for every neighborhood U of 0 in K, there exists a neighborhood V of \(e\) in \(\widehat{\mathbf{G}}\) such that \(\overline{\beta}^{-1}(\mathrm{V}) \subset \mathrm{U}\). Let \(x \mapsto |x|\) be an absolute value on K defining the topology of K (AC, VI, §9, no. 1, prop. 1), and let \(x_0 \in \mathrm{K}\) be such that \(\chi(x_{0}) \neq 1\); let us denote \(\eta = |\chi(x_{0}) - 1| > 0\). Let U be a neighborhood of 0 in K. There exists \(\delta > 0\) such that U contains the set of \(y \in K\) such that \(|y| < \delta\). Let V be the set of characters \(\xi \in \widehat{G}\) such that \(|\langle\xi, x\rangle - 1| < \eta\) for every element \(x \in K\) satisfying \(|x| \leqslant |x_{0}|/\delta\). It is a neighborhood of e in \(\widehat{G}\). If \(y \neq 0\) is such that \(\beta(y)\) belongs to V, we therefore have \(|\chi(xy) - 1| < |\chi(x_{0}) - 1|\) for all x such that \(|x| \leqslant |x_{0}|/\delta\). Consequently, we have \(|x_{0}y^{-1}| > |x_{0}|/\delta\), and hence \(|y| < \delta\), so that \(y \in U\).

Since \(\beta\) is a homeomorphism onto its image, its image is closed in \(\widehat{\mathbf{G}}\) (TG, III, p.~22, cor. 2). But on the other hand, the orthogonal of the image of \(\beta\) is the set of elements \(x\) of \(\mathbf{G}\) such that \(\chi(xy) = 1\) for all \(y \in \mathbf{G}\), and is therefore reduced to \(\{0\}\). The image of \(\beta\) is therefore dense in \(\widehat{\mathbf{G}}\) (corollary 3 of II, p.~228). We conclude that \(\beta\) is surjective.

COROLLARY 1. --- Let K be a locally compact, non-discrete, not necessarily commutative field, and \(\chi\) a nontrivial unitary character of the additive group of K. Let E be a finite-dimensional topological vector space over K. The map \(\theta\) from \(E \times E'\) into U defined by \(\theta(x,\lambda)=\chi(\langle\lambda,x\rangle)\) for \((\lambda,x)\in\mathrm{E}'\times\mathrm{E}\) puts the topological groups E and E' in duality.

Let n be the dimension of E and \((e_{1},\ldots,e_{n})\) a basis of E. It allows one to identify E and \(K^{n}\) (EVT, I, p.~14, th. 2). The result then follows from prop. 19 and prop. 2 of II, p.~206.

We denote by T the group R/Z.

COROLLARY 2. --- a) The group R is in duality with itself with respect to the mapping \((x,y)\mapsto\exp(2i\pi xy)\) , and the dual measure of Lebesgue measure is Lebesgue measure;

\begin{enumerate}
\def\labelenumi{\alph{enumi})}
\setcounter{enumi}{1}
\tightlist
\item
  The groups Z and T are in duality with respect to the map obtained by passing to the quotient from the map from \(Z \times R\) to U such that \((n, x) \mapsto \exp(2i\pi nx)\) . The dual Haar measure of the counting measure on Z is the normalized Haar measure on R/Z.
\end{enumerate}

The group \(\mathbf{R}\) is in duality with itself relative to the mapping \((x,y)\mapsto \exp (2i\pi xy)\) by prop. 19. Let us identify \(\widehat{\mathbf{R}}\) with \(\mathbf{R}\) . The orthogonal of \(\mathbf{Z}\) in \(\widehat{\mathbf{R}} = \mathbf{R}\) is then \(\mathbf{Z}\) , and \(b)\) follows from th. 4 of II, p.~226.

Let \(\alpha\) be the counting measure on Z and \(\gamma\) the normalized Haar measure on T. If \(\beta\) denotes Lebesgue measure on \(\mathbf{R}\), we have \(\gamma = \beta/\alpha\) since these two Haar measures on R/Z have mass 1. The Haar measure \(\widehat{\alpha}\) on \(\widehat{Z} = T\) is the normalized Haar measure (prop. 18), and the Haar measure \(\widehat{\gamma}\) is the counting measure on \(\mathbf{Z}\) (loc. cit.). By Prop. 16 of II, p.~231, the dual measure of \(\beta\) is therefore the measure \(\beta\) .

Remark. --- In particular, one recovers the determination of \(\mathsf{X}(\mathsf{L}^{1}(\mathbf{Z}))\) as done in Example 4 of I, p.~36.

For every integer \(n \geqslant 0\) and \((x, y) \in \mathbf{R}^n \times \mathbf{R}^n\) , we denote by

\[
x \cdot y = \sum_ {j = 1} ^ {n} x _ {j} y _ {j}.
\]

COROLLARY 3. --- Let \(n \geqslant 1\) an integer. The group \(\mathbf{R}^{n}\) is in duality with itself relative to the mapping \((x,y) \mapsto \exp(2i\pi x \cdot y)\) and the dual measure of the Lebesgue measure on \(\mathbf{R}^{n}\) is the Lebesgue measure. The groups \(\mathbf{Z}^{n}\) and \(\mathbf{T}^{n} = \mathbf{R}^{n}/\mathbf{Z}^{n}\) are in duality with respect to the map obtained by passing to the quotient from the map \((n,x) \mapsto \exp(2i\pi x \cdot y)\) , and the dual Haar measure of the counting measure on \(\mathbf{Z}^{n}\) is the normalized Haar measure on \((\mathbf{R}/\mathbf{Z})^{n}\) .

This follows from Lemma 7 of II, p.~225, from Proposition 10 of II, p.~217, and from Corollary 2.

Remark. --- Given a subgroup H of \(R^{n}\) , there corresponds its orthogonal \(H^{\perp}\) , a subgroup of \(\widehat{R^{n}} = R^{n}\) , which is none other than the subgroup associated with H defined in TG, VII, p.~6, no. 3.

In what follows, we shall identify the dual of \(R^{n}\) (resp. of \(T^{n}\) ) with \(R^{n}\) (resp. with \(Z^{n}\) ) by means of the duality of the corollary. In particular, for \(f \in \mathrm{L}^{1}(\mathbf{R}^{n})\) , its Fourier transform is identified with the function on \(R^{n}\) with values in C which, to \(y \in R^{n}\) , associates

\[
\mathcal {F} (f) (y) = \int_ {\mathbf {R} ^ {n}} f (x) \exp (- 2 i \pi x \cdot y) d x.
\]

COROLLARY 4. --- The group \(\mathbf{R}^{*}\) is in duality with the group \(\{-1,1\} \times \mathbf{R}\) by the map \((x,(\sigma,t)) \mapsto \sigma(x/|x|)|x|^{it}\) . The group \(\mathbf{R}_{+}^{*}\) is in duality with \(\mathbf{R}\) by the map \((x,t) \mapsto x^{it}\) .

Indeed, the map \(x \mapsto (x/|x|, \log(|x|))\) is a topological group isomorphism from \(R^{*}\) onto \(\{-1, 1\} \times R\) . The assertion then follows from lemma 7 of II, p.~225, from corollary 2, and from the fact that the unitary characters of \(\{-1, 1\}\) are 1 and \(x \mapsto x\) .

Let \(p\) be a prime number. The field \(\mathbf{Q}_{p}\) of \(p\)-adic numbers is the completion of \(\mathbf{Q}\) with respect to the \(p\)-adic valuation (INT, VII, § 1, no. 6, example, and AC, VI, § 3, no. 4, example 4). For every \(x \in \mathbf{Q}_{p}\), there exists a unique integer \(\nu\geqslant0\) and a unique integer q satisfying \(0\leqslant q<p^{\nu}\) such that \(qp^{-\nu}-x\in Z_{p}\) (A, VII, p.~10, th. 2, applied to the principal ideal domain \(Z_{p}\) and to the set \(R_{p}\) of integers j such that \(0\leqslant j<p\)). We denote by \(\lambda(x)=qp^{-\nu}\) .

PROPOSITION 20. — The mapping \(x \mapsto \exp(2i\pi\lambda(x))\) is a unitary character of \(\mathbf{Q}_{p}\) whose kernel is \(\mathbf{Z}_{p}\) .

For \(x_1\) and \(x_2\) in \(\mathbf{Q}_p\), by definition \(\lambda(x_1 + x_2) - \lambda(x_1) - \lambda(x_2) \in \mathbf{Z}_p \cap \mathbf{Q} = \mathbf{Z}\). The map \(\lambda\) is moreover locally constant since \(\lambda(x + y) = \lambda(x)\) if \(y \in \mathbf{Z}_p\). It then follows that \(x \mapsto \exp(2i\pi\lambda(x))\) is a unitary character of \(\mathbf{Q}_p\) . Since \(\lambda(x) \in \mathbf{Z}\) if and only if \(x \in \mathbf{Z}_p\) , the kernel of this character is \(\mathbf{Z}_p\) .

Recall that one calls normalized Haar measure on the additive group of \(Q_{p}\) the unique Haar measure \(\mu\) such that \(\mu(\mathbf{Z}_{p}) = 1\) (INT, VII, §1, no. 6, example).

COROLLARY. --- a) The group \(\mathbf{Q}_{p}\) is in duality with itself relative to the mapping \((x,y)\mapsto\exp(2i\pi\lambda(xy))\). The normalized Haar measure on \(\mathbf{Q}_{p}\) is then its own dual;

\begin{enumerate}
\def\labelenumi{\alph{enumi})}
\setcounter{enumi}{1}
\tightlist
\item
  The groups \(Z_{p}\) and \(Q_{p}/Z_{p}\) are in duality with respect to the map obtained by passing to the quotient from the map defined by \((z,x)\mapsto\exp(2i\pi\lambda(zx))\) , and the dual measure of the normalized Haar measure on \(Z_{p}\) is the counting measure on \(Q_{p}/Z_{p}\) .
\end{enumerate}

The proof follows step by step that of Cor. 2 of Prop. 19.

\subsection{Euclidean Fourier transform and Fourier series}

* Let \(n \in \mathbf{N}\) . We identify \(\mathbf{R}^n\) and its dual as in Cor. 3 of II, p.~236. The dual measure of the Lebesgue measure is then the Lebesgue measure. We equip \(\mathbf{R}^n\) with the Euclidean norm. For every multi-index \(\alpha \in \mathbf{N}^n\) , and every \(x = (x_1, \ldots, x_n) \in \mathbf{R}^n\) , we shall denote by \(x^\alpha = x_1^{\alpha_1} \cdots x_n^{\alpha_n}\) , and we denote by \(\mathrm{X}^\alpha\) the function \(x \mapsto x^\alpha\) on \(\mathbf{R}^n\) .

Let \(m \in R^{m}\) . Every continuous homomorphism of commutative groups from \(R^{m}\) into \(R^{n}\) is a linear map \(\sigma \in \mathcal{L}(\mathbf{R}^{n}, \mathbf{R}^{m})\) (TG, VII, p.~11, prop. 1). The dual homomorphism \(\widehat{\sigma}\) is identified with the linear map \(^{t}\sigma\) .

The Fourier transform in \(\mathbf{R}^n\) takes on a particularly convenient form in the context of the Schwartz function space and its dual (IV, to appear). We summarize here the main results.

Let \(\mathcal{S}(\mathbf{R}^{n})\) be the space of infinitely differentiable functions \(\varphi\) on \(\mathbf{R}^{n}\) with complex values, such that, for every multi-index \(\alpha\in\mathbf{N}^{n}\) and every integer \(k\in\mathbf{N}\) , the function

\[
x \mapsto \| x \| ^ {k} \partial^ {\alpha} \varphi (x)
\]

is bounded on \(\mathbf{R}^{n}\) . We endow \(\mathcal{S}(\mathbf{R}^{n})\) with the locally convex topology defined by the seminorms

\[
p _ {k, \alpha} \colon \varphi \mapsto \sup _ {x \in \mathbf {R} ^ {n}} \| x \| ^ {k} | \partial^ {\alpha} \varphi (x) |.
\]

We say that \(\mathcal{S}(\mathbf{R}^{n})\) is the space of Schwartz functions on \(\mathbf{R}^{n}\) .

For every \(\alpha \in \mathbf{N}^n\) , the maps \(\varphi \mapsto \partial^\alpha \varphi\) and \(\varphi \mapsto \mathrm{X}^\alpha \varphi\) are continuous from \(\mathcal{S}(\mathbf{R}^n)\) into itself. The space \(\mathcal{S}(\mathbf{R}^n)\) is a topological algebra; it is a Fréchet space and a Montel space (EVT IV, p.~18, def. 4). For every \(p \in [1, +\infty]\) , the space \(\mathcal{S}(\mathbf{R}^n)\) is contained in \(\mathcal{L}^p (\mathbf{R}^n)\) and the canonical injection of \(\mathcal{S}(\mathbf{R}^n)\) into \(\mathcal{L}^p (\mathbf{R}^n)\) is continuous. The image of \(\mathcal{S}(\mathbf{R}^n)\) into \(\mathrm{L}^p (\mathbf{R}^n)\) is dense if \(p \neq +\infty\) .

Since every Schwartz function \(\varphi\) is integrable on \(\mathbf{R}^n\) it admits a Fourier transform denoted \(\widehat{\varphi}\) which is identified with the continuous function on \(\mathbf{R}^n\) defined by

\[
y \mapsto \int_ {\mathbf {R} ^ {n}} \varphi (x) \exp (- 2 i \pi x \cdot y) d x.
\]

The Fourier cotransform of \(\varphi\) is identified, in turn, with the continuous function defined by

\[
y \mapsto \int_ {\mathbf {R} ^ {n}} \varphi (x) \exp (2 i \pi x \cdot y) d x.
\]

Let \(\varphi\in\mathcal{S}(\mathbf{R}^{n})\) . Let \(\alpha\in\mathbf{N}^{n}\) be a multi-index. We have

\[
\mathcal {F} (\partial^ {\alpha} \varphi) = (2 i \pi) ^ {| \alpha |} \mathrm{X} ^ {\alpha} \mathcal {F} (\varphi),
\]

\[
\mathcal {F} (\mathrm{X} ^ {\alpha} \varphi) = (- 2 i \pi) ^ {- | \alpha |} \partial^ {\alpha} (\mathcal {F} (\varphi)).
\]

Proposition 21. — The restriction to \(\mathcal{S}(\mathbf{R}^{n})\) of the Fourier transform is an automorphism of topological vector spaces whose inverse is the restriction of the Fourier cotransform.

Let \(\Lambda \subset \mathbf{R}^n\) be a lattice (TG, VII, p.~4), and let \(\Lambda^{*} \subset \mathbf{R}^{n}\) be the associated lattice (TG, VII, p.~6), also sometimes called the dual lattice.

We call the measure of \(\mathbf{R}^n/\Lambda\) for the Haar measure induced by the Lebesgue measure on \(\mathbf{R}^n\) the covolume of the lattice \(\Lambda\), and denote it by \(V(\Lambda)\) (cf. INT, VIII, §5, no. 5, example). For every function \(f \in \mathcal{S}(\mathbf{R}^n)\) and every \(y \in \mathbf{R}^n\) we have the Poisson formula

\[
\sum_ {x \in \Lambda} f (x + y) = \frac {1}{\mathrm{V} (\Lambda)} \sum_ {z \in \Lambda^ {*}} \widehat {f} (z) \exp (2 i \pi y \cdot z).
\]

Remarks. --- 1) More generally, by the corollary of Proposition 15 of II, p.~230, this formula holds for any complex function integrable on \(\mathbf{R}^n\) such that

\[
\sum_ {x \in \Lambda} | f (x + y) | <   + \infty ,
\]

for all \(y \in R^{n}\) and such that the function on \(T^{n}\) defined by

\[
y \mapsto \sum_ {x \in \Lambda} f (x + y)
\]

is continuous and admits an absolutely convergent Fourier series (cf. below).

\begin{enumerate}
\def\labelenumi{\arabic{enumi})}
\setcounter{enumi}{1}
\tightlist
\item
  There exist functions \(f \in \mathrm{B}(\mathbf{R})\) such that the series \(\sum_{n \in \mathbf{Z}} f(n)\) diverge (exercise 4 of II, p.~263).
\end{enumerate}

Example. --- Let Q be a positive definite quadratic form on \(R^{n}\). The function defined by \(\varphi(x)=\exp(-\pi\mathrm{Q}(x))\) belongs to \(\mathcal{S}(\mathbf{R}^{n})\) . There exists \(\sigma\in\mathrm{GL}(n,\mathbf{R})\) such that \(\mathrm{Q}(x)=\|\sigma(x)\|^{2}\) for all \(x\in R^{n}\). The Fourier transform of \(\varphi\) is given for all \(y\in R^{n}\) by

\[
\widehat {\varphi} (y) = \frac {1}{| \det (\sigma) |} \exp (- \pi \mathrm{Q} ^ {*} (y))
\]

where \(\mathrm{Q}^{*}(y)=\|^{t}\sigma^{-1}(y)\|^{2}\) (cf. INT, IX, §6, nos. 4–5 and exercise 1, c) of II, p. 262).

DEFINITION 6. --- We call the space of tempered distributions on \(\mathbf{R}^{n}\) the dual space of \(\mathcal{S}(\mathbf{R}^{n})\) endowed with the topology of bounded convergence. It is denoted \(\mathcal{S}'(\mathbf{R}^{n})\) .

Since \(\mathcal{S}(\mathbf{R}^n)\) is bornological, the space \(\mathcal{S}'(\mathbf{R}^n)\) is complete and bornological (EVT, III, p.~24, cor. 1 and 2). Since \(\mathcal{S}(\mathbf{R}^n)\) is a Montel space, so is \(\mathcal{S}'(\mathbf{R}^n)\) (EVT, IV, p.~19, prop. 9).

Let \(\alpha\in\mathbf{N}^{n}\) . We still denote by \(f\mapsto\mathrm{X}^{\alpha}f\) the transpose of the endomorphism \(\varphi\mapsto\mathrm{X}^{\alpha}\varphi\) of \(\mathcal{S}(\mathbf{R}^{n})\) , and we denote by \(f\mapsto\partial^{\alpha}f\) the endomorphism of \(\mathcal{S}'(\mathbf{R}^{n})\) defined by

\[
\langle \partial^ {\alpha} f, \varphi \rangle = (- 1) ^ {| \alpha |} \langle f, \partial^ {\alpha} \varphi \rangle
\]

for \(f\in \mathcal{S}'(\mathbf{R}^n)\) and \(\varphi \in \mathcal{S}(\mathbf{R}^n)\) .

Let \(f\) be a linear map from \(\mathcal{S}(\mathbf{R}^{n})\) into \(\mathbf{C}\) . Then \(f\) is a tempered distribution if, and only if, for every family \((\mathrm{M}_{k,\alpha})_{(k,\alpha)\in\mathbf{N}\times\mathbf{N}^{n}}\) in \(\mathbf{R}_{+}\), the linear form \(f\) is bounded on the set of functions \(\varphi\in\mathcal{S}(\mathbf{R}^{n})\) such that for all \((k,\alpha)\in\mathbf{N}\times\mathbf{N}^{n}\) , one has \(p_{k,\alpha}(\varphi)\leqslant\mathrm{M}_{k,\alpha}\) .

A sequence \((f_{m})_{m\in\mathbf{N}}\) of tempered distributions converges to a tempered distribution f if, and only if, one has \(\langle f_{m},\varphi\rangle\to\langle f,\varphi\rangle\) for all \(\varphi\in\mathcal{S}(\mathbf{R}^{n})\) .

Example. --- A measure \(\nu\) on \(\mathbf{R}^n\) is said to be tempered if there exists a positive integer \(r\) such that the continuous map \(x \mapsto (1 + \|x\|)^{-r}\) is \(\nu\)-integrable on \(\mathbf{R}^n\). The restriction of \(\nu\) to \(\mathcal{S}(\mathbf{R}^n)\) is a tempered distribution. It is zero if and only if the measure \(\nu\) is zero.

Let \(p \in [1, +\infty]\) and \(f \in \mathcal{L}^{p}(\mathbf{R}^{n})\). Then the measure \(f \cdot dx\) of density f with respect to Lebesgue measure is tempered. In particular, Lebesgue measure \(\mu\) on \(R^{n}\) is tempered, and every bounded measure on \(R^{n}\) is tempered.

For every \(p \in [1, +\infty]\) one can identify \(\mathrm{L}^{p}(\mathbf{R}^{n})\) with a subspace of \(\mathcal{S}'(\mathbf{R}^{n})\) by the linear map \(f \mapsto f \cdot dx\) ; this map is continuous.

DEFINITION 7. --- We call the transpose of the Fourier transform on \(\mathcal{S}(\mathbf{R}^n)\) (respectively, of the Fourier cotransform) the Fourier transform (respectively, Fourier cotransform) on \(\mathcal{S}'(\mathbf{R}^n)\), and denote them by \(\mathcal{F}\) and \(\overline{\mathcal{F}}\), respectively.

For \(f \in \mathcal{S}'(\mathbf{R}^n)\) , the tempered distribution \(\mathcal{F}(f)\) (resp. \(\overline{\mathcal{F}}(f)\) ) is defined by \(\varphi \mapsto \langle f, \mathcal{F}(\varphi) \rangle\) for \(\varphi \in \mathcal{S}(\mathbf{R}^n)\) (resp. by \(\varphi \mapsto \langle f, \overline{\mathcal{F}}(\varphi) \rangle\) ).

The Fourier transform on \(\mathcal{S}'(\mathbf{R}^n)\) is an automorphism of topological vector spaces whose inverse is the Fourier cotransform \(\overline{\mathcal{F}}\) .

PROPOSITION 22. --- Let f be a tempered distribution belonging to \(\mathcal{M}^{1}(\mathbf{R}^{n})\) (respectively to \(\mathrm{L}^{2}(\mathbf{R}^{n})\) ). The Fourier transform of \(f\) in \(\mathcal{S}'(\mathbf{R}^{n})\) is the tempered distribution associated with the Fourier transform of f in \(\mathcal{C}_{0}(\mathbf{R}^{n})\) (respectively in \(L^{2}(\mathbf{R}^{n})\) ). The same holds for the Fourier cotransform.

Remark. --- The elementary formulas concerning the Fourier transform of measures remain valid for the Fourier transform of tempered distributions.

Thus, if \(\alpha \in \mathbf{N}^n\) and \(f\in \mathcal{S}'(\mathbf{R}^n)\) , one has

\[
\mathcal {F} (\partial^ {\alpha} f) = (2 i \pi) ^ {| \alpha |} X ^ {\alpha} \mathcal {F} (f),
\]

\[
\mathcal {F} (X ^ {\alpha} f) = (- 2 i \pi) ^ {- | \alpha |} \partial^ {\alpha} (\mathcal {F} (f)).
\]

Let \(y \in \mathbf{R}^n\) . Denote by \(\gamma(y)\) the endomorphism of \(\mathcal{S}'(\mathbf{R}^n)\) defined by

\[
\langle \boldsymbol {\gamma} (y) f, \varphi \rangle = \langle f, \boldsymbol {\gamma} (- y) \varphi \rangle
\]

for \(f \in \mathcal{S}'(\mathbf{R}^n)\) and \(\varphi \in \mathcal{S}(\mathbf{R}^n)\) . Denote by \(e_y\) the character of \(\mathbf{R}^n\) such that \(e_y(x) = \exp(2i\pi x \cdot y)\) . Then \(e_y \in \mathcal{S}'(\mathbf{R}^n)\) . We have \(\mathcal{F}(e_y) = \varepsilon_y\) and more generally

\[
\mathcal {F} (e _ {y} f) = \boldsymbol {\gamma} (y) \mathcal {F} (f)
\]

for all \(f \in \mathcal{S}'(\mathbf{R}^n)\) . *

Let \(n \geqslant 1\) be an integer and \(G = T^{n}\) , equipped with the normalized Haar measure. The dual group of G is identified with \(Z^{n}\) by the map \(h \mapsto \chi_{h}\) , where \(\chi_{h}\) is the unitary character of \(T^{n}\) obtained by passing to the quotient from the character \(x \mapsto \exp(2i\pi h \cdot x)\) of \(R^{n}\) (corollary 3 of II, p.~236). The Fourier transform of a measure \(\mu\) on \(T^{n}\) is identified with the family \((\widehat{\mu}(h))_{h \in \mathbf{Z}^{n}}\) where

\[
\widehat {\mu} (h) = \int_ {\mathbf {T} ^ {n}} e ^ {- 2 i \pi h \cdot x} d \mu (x).
\]

The series

\[
\sum_ {h \in \mathbf {Z} ^ {n}} \widehat {\mu} (h) \chi_ {h}
\]

is called the Fourier series of μ.

If \(f \in L^{1}(\mathbf{T}^{n})\) is such that its Fourier series converges absolutely in \(L^{1}(\mathbf{Z}^{n})\) , we then have \(f \in \mathcal{C}(\mathbf{T}^{n})\) and

\[
f (x) = \sum_ {h \in {\bf Z} ^ {n}} \widehat {f} (h) e ^ {2 i \pi h \cdot x}
\]

for all \(x \in \mathbf{T}^n\) (Theorem 3 of II, p.~222), where

\[
\widehat {f} (h) = \int_ {\mathbf {T} ^ {n}} f (x) e ^ {- 2 i \pi h \cdot x} d x, \qquad h \in \mathbf {Z} ^ {n}.
\]

For \(f \in \mathrm{L}^2(\mathbf{T}^n)\) , the Fourier inversion formula (prop. 12 of II, p.~219) states that, if the series with general term \(\widehat{f}(h)\) converges absolutely, we have

\[
f (x) = \sum_ {h \in \mathbf {Z} ^ {n}} \widehat {f} (h) e ^ {2 i \pi h \cdot x}
\]

for almost every x in \(T^{n}\) .

However, even if f is continuous, the Fourier series of f does not in general converge to \(f(x)\) for all x (exerc. 30 of II, p.~274). The following result is all the more useful.

PROPOSITION 23 (Fejér's Theorem). --- Let \(n \geqslant 1\) be an integer. For every \(h = (h_i) \in \mathbf{Z}^n\) , we denote by \(|h| = \sup_i |h_i|\) . Let \(f \in \mathcal{C}(\mathbf{T}^n)\) . For every integer \(N \geqslant 1\) , let \(f_N\) be the function on \(\mathbf{T}^n\) such that

\[
f_{\mathrm{N}}(x) = \sum_{\substack{h\in \mathbf{Z}^{n}\\ |h|\leqslant \mathrm{N}}}\widehat{f} (h)e^{2i\pi h\cdot x}\prod_{j = 1}^{n}\Bigl (1 - \frac{|h_{j}|}{\mathrm{N}}\Bigr)
\]

for \(x\in \mathbf{T}^n\). Then \(f_{\mathrm{N}}\) converges to \(f\) in \(\mathcal{C}(\mathbf{T}^n)\).

Lemma 11. --- For every N ≥ 1, let \(\mu_{\mathrm{N}}\) be the measure on \(\mathbf{T}^{n}\) whose density is the continuous map

\[
\mathrm{F}_{\mathrm{N}}\colon x\mapsto \sum_{\substack{h\in \mathbf{Z}\\ |h|\leqslant \mathrm{N}}}e^{2i\pi h\cdot x}\prod_{j = 1}^{n}\Bigl (1 - \frac{|h_{j}|}{\mathrm{N}}\Bigr).
\]

The sequence of measures \((\mu_{\mathrm{N}})_{\mathrm{N}\geqslant1}\) converges to \(\varepsilon_{0}\) in the space \(\mathcal{M}^{1}(\mathbf{T}^{n})\) endowed with the topology of compact convergence in \(\mathcal{C}(\mathbf{T}^{n})\) .

We reduce to the case n = 1 by noting that \(\mu_{N}\) is the product of measures of the same type for n = 1. It then suffices to verify that the sequence ( \(\mu_{N}\) ) satisfies the hypotheses of Lemma 4 of INT, VIII, §2, n°7 with a = 0.

To this end, we first note that \(F_{N}\) is the Fourier cotransform of the map \(\varphi_{\mathrm{N}}\colon h\mapsto(1-|h|/\mathrm{N})\) on Z. This is written as \(\varphi_{N}=N^{-1}\psi_{N}*\widetilde{\psi}_{N}\) , where \(\psi_{N}\) is the characteristic function of the set defined by \(-N/2<|h|\leqslant N/2\) . Consequently, \(F_{N}=N^{-1}|\overline{\mathcal{F}}(\psi_{N})|^{2}\geqslant0\). Thus, \(\mu_{N}\) is a positive measure; we have \(\mu_{\mathrm{N}}(\mathbf{T})=1\), which proves (i) and (iii) in loc. cit.

Let us prove condition (ii) of loc. cit. Let U be an open neighborhood of 0 in T. It suffices to show that \(\mu_{\mathrm{N}}(\mathrm{U}) \to 1\) when \(\mathrm{N} \to +\infty\). Let K be a compact symmetric neighborhood of 0 such that \(\mathrm{K}^2 \subset \mathrm{U}\) and \(\psi\) the characteristic function of K. Set \(\varphi = \psi * \psi\) . This is an element of A(T) with support contained in U. The real number \(m = \varphi(0)\) is the measure of the set K, and hence \(m > 0\). Moreover, we have \(0 \leqslant \varphi \leqslant m\) since \(\varphi(x)\) is the measure of the set K \(\cap x\) K. We have

\[
\mu_ {\mathrm{N}} (\mathrm{U}) \geqslant \frac {1}{m} \int_ {\mathbf {T}} \varphi (x) \mu_ {\mathrm{N}} (x) = \frac {1}{m} \sum_ {h \in \mathbf {Z}} \mathcal {F} (\varphi) (h) \varphi_ {\mathrm{N}} (h)
\]

by the transposition properties of the Fourier transform (prop. 13 of II, p.~221). Since \(\varphi \in \mathrm{A}(\mathbf{T})\) , its Fourier transform belongs to \(\mathrm{L}^1 (\mathbf{Z})\) and \(\varphi\) satisfies the Fourier inversion formula (prop. 11 of II, p.~217). Since \(\varphi_{\mathrm{N}}(h)\to 1\) for all \(h\in \mathbf{Z}\) and \(|\varphi_{\mathrm{N}}(h)|\leqslant 1\) , Lebesgue's theorem (INT, IV, §3, no. 7, th. 6) and the Fourier inversion formula imply that

\[
\begin{array}{r l} & {\underset {\mathrm{N} \to + \infty} {\liminf} \mu_ {\mathrm{N}} (\mathrm{U}) \geqslant \frac {1}{m} \underset {\mathrm{N} \to + \infty} {\lim} \sum_ {h \in \mathbf {Z}} \mathcal {F} (\varphi) (h) \varphi_ {\mathrm{N}} (h) =} \\ & {\qquad \qquad \qquad \qquad \frac {1}{m} \sum_ {h \in \mathbf {Z}} \mathcal {F} (\varphi) (h) = \frac {1}{m} \varphi (0) = 1.} \end{array}
\]

Let us prove the proposition. We have \(f * \mathrm{F}_{\mathrm{N}} = f_{\mathrm{N}}\) for \(\mathrm{N} \geqslant 1\) . The regular representation \(\gamma\) of \(\mathbf{T}^n\) into \(\mathcal{C}(\mathbf{T}^n)\) (INT, VIII, §2, no. 3) is continuous and satisfies \(f * \mathrm{F}_{\mathrm{N}} = \gamma(\mu_{\mathrm{N}})f\) (INT, VIII, §4, no. 5, prop. 5 (iv)). The map \(\mu \mapsto \gamma(\mu)f\) is continuous from \(\mathcal{M}^1(\mathbf{T}^n)\) into \(\mathcal{C}(\mathbf{T}^n)\) (INT, VI, §1, no. 6, prop. 14). By the lemma, we therefore have

\[
\lim _ {N \to + \infty} f _ {N} = \lim _ {N \to + \infty} f * F _ {N} = \lim _ {N \to + \infty} \gamma (\mu_ {N}) f = \gamma (\varepsilon_ {0}) (f) = f
\]

in \(\mathcal{C}(\mathbf{T}^n)\) .

Note. --- There exist functions \(f \in \mathrm{L}^{1}(\mathbf{T})\) whose Fourier series diverges at every point \(x \in T\) (Kolmogorov's theorem, cf.~Exercise 51 of II, p.~289).

A theorem of Carleson \(^{(1)}\) proves that the symmetric partial sums of the Fourier series of f converge to \(f(x)\) for almost all \(x \in T\) if \(f \in \mathcal{L}^{2}(T)\) .

\section{CLASSIFICATION}

\subsection{Groups generated by a compact subset}

Lemma 1. --- Let H be a locally compact group, and let R be one of the groups R or Z. Let \(\varphi\) be a continuous morphism from R to H. If \(\varphi\) is not a topological isomorphism of R onto a subgroup of H, then the image of R in H is relatively compact.

Let I be the image of \(\varphi\) . Replacing H by \(\bar{I}\) , we may assume that I is dense in H. We must then show that H is compact if \(\varphi\) is not a topological isomorphism of R onto I.

Suppose there exists a neighborhood V of \(e\) in H and an integer \(M > 0\) such that, for all \(t > \mathrm{M}\) in R, we have \(\varphi(t) \notin \mathrm{V}\). Then \(\varphi\) is injective: if \(\varphi(u) = e\), one has \(\varphi(nu) = e\) for every integer \(n \geqslant 1\), hence the set \(\mathbf{N}u\) is bounded in R, which means that \(u = 0\). The restriction of \(\varphi\) to \([-M, M] \cap R\) is therefore a homeomorphism onto its image, which contains \(V \cap I\). The restriction of \(\varphi^{-1}\) to \(V \cap I\) being continuous, it follows that \(\varphi\) is a topological isomorphism of R onto I.

Suppose now that \(\varphi\) is not a topological isomorphism from R onto I. Let W be a relatively compact open neighborhood of \(e\) in H, and V a symmetric neighborhood of \(e\) such that \(\mathrm{V}^2 \subset \mathrm{W}\). For every \(x \in \mathrm{H} = \overline{\mathrm{I}}\), there exists an element \(s \in \mathrm{R}\) such that \(x \in \varphi(s)\mathrm{V}\). From the preceding paragraph and the hypothesis on \(\varphi\), there exists \(t \in \mathrm{R}\) such that \(t > |s|\) and \(\varphi(t) \in \mathrm{V}\). We then have \(x \in \varphi(t + s)\varphi(t)^{-1}\mathrm{V} \subset \varphi(t + s)\mathrm{W}\), and \(t + s > 0\). Consequently, the open sets \(\varphi(u)\mathrm{W}\) for \(u > 0\) form an open covering of H. Since W is relatively compact, there exists an integer \(n \geqslant 1\) and elements \(u_1, \ldots, u_n\) of R, strictly positive, such that \(\overline{\mathrm{W}} \subset \bigcup_{1 \leqslant i \leqslant n} \varphi(u_i)\mathrm{W}\). Let U be the largest of the \(u_i\) .

Let \(x \in \mathrm{H}\) and let \(s = \inf \{t \in \mathrm{R} | t \geqslant 0, \varphi(t)x^{-1} \in \overline{\mathrm{W}}\}\) . Since \(\overline{\mathrm{W}}\) is compact, we then have \(\varphi(s)x^{-1} \in \overline{\mathrm{W}}\) . There exists an integer \(i\) such that \(\varphi(s)x^{-1} \in \varphi(u_i)\mathrm{W}\) , whence \(\varphi(s - u_i)x^{-1} \in \overline{\mathrm{W}}\) . The definition of \(s\) implies \(s - u_i < 0\) , whence \(s \leqslant \mathrm{U}\). It follows that \(\mathrm{H} = \varphi([0,\mathrm{U}] \cap \mathrm{R})\overline{\mathrm{W}}\) is compact.

Lemma 2. --- If G is generated by a compact subset V, then there exists an integer \(n \geqslant 0\) and a discrete subgroup D of G isomorphic to \(Z^{n}\) such that G/D is compact.

Up to replacing V by \(V \cup V^{-1}\) , we may assume that V is symmetric; the hypothesis then means that G is the union of the sets \(V^{n}\) where \(n \in N\) .

Since \(V^{2}\) is compact, there exists an integer \(k \geqslant 1\) and elements \(x_{1}, \ldots, x_{k} \in G\) such that \(V^{2} \subset \bigcup_{1 \leqslant i \leqslant k} x_{i}V\) . Let \(D_{0}\) be the subgroup of G generated by the family \((x_{i})_{1 \leqslant i \leqslant k}\). We have \(V^{2} \subset D_{0}V\), whence by induction \(V^{n} \subset D_{0}V\) for every integer \(n \geqslant 1\), and hence \(G = D_{0}V\) since V generates G. Let then J be a subset of \(\{1, 2, \ldots, k\}\) such that the subgroup D generated by the family \((x_{i})_{i \in J}\) is topologically isomorphic to \(\mathbf{Z}^{\mathrm{Card(J)}}\) and maximal for this property. Let us show that G/D is compact.

Let \(p\) be the canonical surjection from G onto G/D. Let \(i \in \{1, 2, \ldots, k\} - J\) . If the subgroup \(H_i\) of G/D generated by \(p(x_i)\) is topologically isomorphic to \(Z\) , the subgroup of G generated by D and \(x_i\) is discrete and the map \((d, n) \mapsto dx_i^n\) is an isomorphism from D × \(Z\) onto this subgroup, contrary to the maximality of J. Lemma 1 therefore implies that \(\overline{H}_i\) is compact. Hence G/D = \((\prod_{i \not\in J} \overline{H}_i)p(V)\) is compact.

Lemma 3. — Let A and B be commutative groups such that A is divisible. Let C be a subgroup of B and \(\varphi\) a morphism from C to A. There exists a morphism from B to A extending \(\varphi\) .

Let \(\mathcal{O}\) be the set of pairs \((\mathrm{X},f)\), where \(\mathrm{X}\) is a subgroup of B containing C and \(f\) a morphism from \(\mathrm{X}\) to A extending \(\varphi\). Order \(\mathcal{O}\) by the relation " \(\mathrm{X} \subset \mathrm{X}'\) and \(f'\) extends \(f\) ". One verifies that \(\mathcal{O}\) is inductive. Let \((\mathrm{X},f)\) be a maximal element of \(\mathcal{O}\) (E, III, p.~20, th. 2). If \(\mathrm{X} \neq \mathrm{B}\) , take an element \(b\) of B-X and let \(\mathrm{X}'\) be the subgroup generated by \(\mathrm{X}\) and \(b\) . The commutativity of B shows that \(\mathrm{X}'\) is the set of elements \(b^n x\) for \(n \in \mathbf{Z}\) and \(x \in \mathrm{X}\) . Suppose first that \(b^n \notin \mathrm{X}\) for every integer \(n \neq 0\) and define \(f'\) as a morphism from \(\mathrm{X}'\) to A by taking an arbitrary element \(y \in \mathrm{A}\) and setting \(f'(b^n x) = y^n f(x)\) for all \(n \in \mathbf{Z}\) and every \(x \in \mathrm{X}\) . Since A is commutative, \(f'\) is a morphism, and it extends \(f\) . Suppose now that there exists \(n \neq 0\) such that \(b^n \in \mathrm{X}\) and let \(m > 0\) such that \(m\mathbf{Z} = \{n \in \mathbf{Z} \mid b^n \in \mathrm{X}\}\) . Since A is divisible, there exists an element \(y \in \mathrm{A}\) such that \(y^m = f(b^m)\) . One then extends \(f\) to a morphism from \(\mathrm{X}'\) to A by \(f'(b^n x) = y^n f(x)\) for \(n \in \{0,1,\ldots,m-1\}\) and \(x \in \mathrm{X}\) . In both cases, \((\mathrm{X},f)\) would not be maximal. Hence we have \(\mathrm{X} = \mathrm{B}\) and the lemma is proved.

Remark. --- In the language of categories, the lemma says that divisible groups are injective objects in the category of commutative groups; cf.~A, VII, p.~53, exerc. 3.

PROPOSITION 1. --- The following conditions are equivalent:

\begin{enumerate}
\def\labelenumi{(\roman{enumi})}
\item
  G is generated by a compact subset;
\item
  there exist positive integers p and q and a compact group K such that G is isomorphic to \(R^{p} \times Z^{q} \times K\) ;
\item
  there exists an integer \(n \geqslant 0\) such that \(\widehat{G}\) is locally isomorphic to \(R^{n}\) ;
\item
  there exist positive integers p and q and a discrete group D such that \(\widehat{G}\) is isomorphic to \(R^{p} \times T^{q} \times D\) .
\item
  \(\Longrightarrow\) (iii): if G has property (i), there exists an integer \(n \geqslant 0\) and a subgroup D of G isomorphic to \(Z^{n}\) such that G/D is compact (lemma 2). Then \(D^{\perp}\) , which is identified with the dual of G/D, is discrete (th. 4 of II, p.~226 and prop. 18 of II, p.~233). Hence \(\widehat{G}\) is locally isomorphic to \(\widehat{G}/D^{\perp}\) , that is to say to \(\widehat{D}\) , which is isomorphic to \(T^{n}\) (th. 4 of II, p.~226 and prop. 18 of II, p.~233). Now \(T^{n}\) is locally isomorphic to \(R^{n}\) .
\item
  \(\Longrightarrow\) (iv): if \(\widehat{\mathbf{G}}\) is locally isomorphic to \(\mathbf{R}^n\) , there exists an integer \(p\) such that \(0 \leqslant p \leqslant n\) so that the neutral component \(\widehat{\mathbf{G}}_0\) of \(\widehat{\mathbf{G}}\) is an open subgroup isomorphic to \(\mathbf{R}^p \times \mathbf{T}^{n-p}\) (TG, VII, p.~13, th. 1). In particular, \(\widehat{\mathbf{G}}_0\) is a divisible group. Applying lemma 3 to the identity map of the subgroup \(\widehat{\mathbf{G}}_0\) of the group \(\widehat{\mathbf{G}}\) into the divisible group \(\widehat{\mathbf{G}}_0\) , there exists a morphism \(\pi\) of \(\widehat{\mathbf{G}}\) into \(\widehat{\mathbf{G}}_0\) which is the identity map on \(\widehat{\mathbf{G}}_0\) . Consequently, we have \(\pi \circ \pi = \pi\) , and \(\pi\) is a projector. It is continuous, since its restriction to the open subgroup \(\widehat{\mathbf{G}}_0\) is. Hence, \(\widehat{\mathbf{G}}\) is the direct product of \(\widehat{\mathbf{G}}_0\) and the subgroup \(\overline{\pi}^{-1}(e)\) , which is discrete since it is isomorphic to \(\widehat{\mathbf{G}} / \widehat{\mathbf{G}}_0\) (TG, III, p.~47, cor.) (iv) \(\Longrightarrow\) (ii): follows from prop. 2 of II, p.~206, from prop. 18 of II, p.~233 and from corollary 3 of II, p.~236.
\item
  \(\Longrightarrow\) (i): for every compact group K, the group \(R^{p} \times Z^{q} \times K\) is generated by the compact set \([0,1]^{p} \times \{0,1\}^{q} \times K\) .
\end{enumerate}

COROLLARY 1. --- Suppose that G is generated by a compact neighborhood of e.

\begin{enumerate}
\def\labelenumi{\alph{enumi})}
\item
  There exists a compact subgroup K of G and positive integers p and q such that G is isomorphic to \(R^{p} \times Z^{q} \times K\) .
\item
  Conversely, let K be a compact group, let p and q be positive integers, and let G be a group isomorphic to \(R^{p} \times Z^{q} \times K\). Then K is the unique maximal compact subgroup of G, and the integers \((p,q)\) are uniquely determined by G.
\end{enumerate}

Assertion \(a\) ) follows from prop. 1. Let then K be a compact group, and \(p\) , \(q\) positive integers. Suppose that G is isomorphic to the group \(\mathbf{R}^{p} \times \mathbf{Z}^{q} \times \mathrm{K}\) , and we identify G with this group. Under the canonical projection of G onto \(\mathbf{R}^{p} \times \mathbf{Z}^{q}\) , the image of every compact subgroup of G is a compact subgroup of \(\mathbf{R}^{p} \times \mathbf{Z}^{q}\); therefore it is reduced to the identity element. Therefore \(\mathrm{K}' \subset \mathrm{K}\) and K is the largest compact subgroup of G. The subgroup \(\mathbf{R}^{p} \times \mathrm{K}\) is also unique because \(\mathbf{R}^{p}\) is the neutral component of G/K. Taking into account TG, VII, p.~13, cor. 3, the integer \(p\) is uniquely determined by G. Since \(\mathrm{G}/(\mathbf{R}^{p} \times \mathrm{K})\) is isomorphic to \(\mathbf{Z}^{q}\) , the integer \(q\) is also uniquely determined by G.

Note. --- By duality, \(\hat{G}\) is isomorphic to \(R^{p} \times T^{q} \times D\) (prop. 3 (iv)), where the subgroups \(R^{p} \times T^{q}\) and \(T^{q}\), and the integers p and q, are uniquely determined.

COROLLARY 2. --- The following conditions are equivalent:

\begin{enumerate}
\def\labelenumi{(\roman{enumi})}
\item
  The groups G and \(\hat{G}\) are generated by compact subsets;
\item
  There exist positive integers n and m such that G is locally isomorphic to \(R^{m}\) and \(\hat{G}\) to \(R^{n}\) ;
\item
  There exist positive integers p, q, and r and a finite group A such that G is isomorphic to \(R^{p} \times T^{q} \times Z^{r} \times A\) ;
\item
  There exist positive integers p, q, and r and a finite group A such that \(\hat{G}\) is isomorphic to a product \(R^{p} \times Z^{q} \times T^{r} \times A\) .
\end{enumerate}

We have (i) \(\Leftrightarrow\) (ii) by Prop. 1, and (iii) \(\Leftrightarrow\) (iv) by duality, hence (iii) \(\Rightarrow\) (i). Finally, if (i) is true, then \(\widehat{\mathbf{G}}\) is isomorphic to \(\mathbf{R}^{p} \times \mathbf{T}^{r} \times \mathbf{D}\), where D is discrete (prop. 1). The group D is generated by a compact subset and is therefore finite. Consequently, there exists \(q \geqslant 0\) such that D is isomorphic to \(\mathbf{Z}^{q} \times \mathbf{A}\) , where A is a finite group (A, VII, p.~22, th. 3).

Note. --- With the notations of Cor. 2, if one identifies G with \(R^{p} \times T^{q} \times Z^{r} \times A\) the subgroup \(R^{p} \times T^{q}\) is the identity component of G, the subgroup \(T^{q} \times A\) is its largest compact subgroup and \(T^{q}\) is the identity component of it; the integers p, q, r are uniquely determined by G according to the preceding remark, and the group A is determined by G up to isomorphism.

PROPOSITION 2. --- Suppose G is compact. There exists a decreasing filtered family \((\mathrm{H}_{i})_{i\in\mathrm{I}}\) of closed subgroups of G such that

\begin{enumerate}
\def\labelenumi{\alph{enumi})}
\item
  the group G is identified with the projective limit of the \(\mathrm{G / H_i}\) ;
\item
  for every i, there exists an integer \(q \geqslant 0\) and a finite group A such that \(G/H_{i}\) is isomorphic to \(T^{q} \times A\) .
\end{enumerate}

Indeed, \(\widehat{\mathbf{G}}\) is discrete (prop. 18 of II, p.~233), hence is the union of an increasing filtered family \((\mathrm{D}_i)_{i\in \mathrm{I}}\) of finitely generated subgroups. Set \(\mathrm{H}_i = \mathrm{D}_i^\perp\); the group G is identified with the projective limit of the \(\mathrm{G / H_i}\) (II, p. 234, cor. 3), and \((\mathrm{H}_i)\) is a decreasing filtered family.

Let \(i \in I\) . There exists an integer \(q \geqslant 0\) and a finite group A such that the group \(D_{i}\) is isomorphic to \(Z^{q} \times A\) (A, VII, p.~22, th. 3), hence \(G/H_{i}\) is isomorphic to \(T^{q} \times \widehat{A}\) (cf.~corollary 1 of II, p.~232).

COROLLARY. --- If the group G is generated by a compact subset, then it is a projective limit of groups isomorphic to groups of the form \(R^{p} \times T^{q} \times Z^{r} \times A\) , where A is a finite group and p, q, r are positive integers.

By (ii) of prop. 1 of II, p.~246, the corollary follows from prop. 2.

\subsection{General case}

In this section, G denotes a locally compact commutative group.

PROPOSITION 3. --- a) There exists an integer \(n \geqslant 0\) and a subgroup L such that G is the direct product of L and a subgroup isomorphic to \(\mathbf{R}^n\) , and moreover L admits a compact open subgroup K such that L/K is discrete;

\begin{enumerate}
\def\labelenumi{\alph{enumi})}
\setcounter{enumi}{1}
\tightlist
\item
  The group G is the union of an increasing filtered family of open subgroups, each being a projective limit of groups isomorphic to groups of the form \(R^{p} \times T^{q} \times Z^{r} \times A\) , where A is a finite group and p, q, r are positive integers.
\end{enumerate}

We prove \(b\) ). For every compact neighborhood V of \(e\) , let \(\mathrm{G_V}\) denote the subgroup of G generated by V. It is open, and by the corollary of prop. 2, the group \(\mathrm{G_V}\) is a projective limit of groups of the form \(\mathbf{R}^p\times \mathbf{T}^q\times \mathbf{Z}^r\times \mathrm{A}\) , where A is a compact group and \(p,q\) and \(r\) in \(\mathbf{N}\) . When V ranges over the compact neighborhoods of \(e\) , these subgroups \(\mathrm{G_V}\) form a filtered family (since \(G_{V}\) and \(G_{W}\) are contained in \(G_{V\cup W}\) for all compact neighborhoods V and W of e). Finally, the group G is the union of the subgroups \(G_{V}\) .

Let H be an open subgroup of G generated by a compact neighborhood of e. There exist a compact group K and positive integers p and q such that H is isomorphic to \(R^{p} \times Z^{q} \times K\) (prop. 1 of II, p.~246); identify H with this product. The canonical surjection of H onto the divisible group \(R^{p}\) extends to a morphism \(\pi\) of G onto \(R^{p}\) (lemma 3 of II, p.~245). It is a projector \(\pi\) of G onto \(R^{p}\), which is continuous since its restriction to the open subgroup H is continuous. Hence G is the direct product of \(R^{p}\) and of the kernel L of \(\pi\) (TG, III, p.~47, cor.). We have \(Z^{q} \times K = H \cap L\) , hence \(Z^{q} \times K\) is an open subgroup of L. Thus K is a compact open subgroup of L, and consequently L/K is discrete.

PROPOSITION 4. --- Let B \(_{G}\) be the set of elements of G that generate a relatively compact subgroup of G. Then B \(_{G}\) is a closed subgroup of G and B \(_{G}^{\perp}\) is the identity component of \(\widehat{G}\) .

The set \(B_{G}\) is a subgroup of G, since the product of two compact subsets of G is a compact subset of G.

Let H be an open subgroup of G generated by a compact neighborhood of e. There exist positive integers p and q and a compact group K such that the subgroup H is isomorphic to \(R^{p} \times Z^{q} \times K\) (prop. 1 of II, p.~246). If one identifies these groups, one sees that \(B_{G} \cap H = K\) is closed in H. Since the family of open subgroups H generated by compact neighborhoods is an open cover of G (for example, x belongs to the subgroup generated by \(U \cap \{x\}\) for every fixed compact neighborhood U of e), it follows that \(B_{G}\) is closed.

Now let us calculate \(B_{G}^{\perp}\). By prop. 3, a), there exists a positive integer \(n \geqslant 0\) and a group L admitting an open compact subgroup such that G can be identified with \(R^{n} \times L\) . Then \(B_{G}\) is identified with \(\{0\} \times B_{L}\) , and \(B_{G}^{\perp}\) with \(R^{n} \times B_{L}^{\perp}\). We are thus reduced to the case where G = L admits a compact open subgroup K.

We then have \(K \subset B_{G}\). If one identifies \(\widehat{G/K}\) with \(K^{\perp}\) (Theorem 4 of II, p.~226), the orthogonal \((\mathrm{B}_{\mathrm{G}}/\mathrm{K})^{\perp}\) of \(B_{G}/K\) in \(\widehat{G/K}\) is identified with \(B_{G}^{\perp}\). But on the other hand \(B_{G}/K = B_{G/K}\), and since \(K^{\perp}\) is an open subgroup of \(\widehat{G}\) (cor. 2 of II, p.~233), the identity component \(\widehat{G}_{0}\) of \(\widehat{G}\) is also the identity component of \(K^{\perp}\). Thus the assertion for the discrete group G/K is equivalent to that for G.

Finally, suppose that G is discrete. The group \(\widehat{G}\) is then compact (prop. 18 of II, p.~233). The neutral component \((\widehat{\mathrm{G}})_{0}\) is the intersection of the open subgroups of \(\widehat{G}\) (TG, III, p.~35, prop. 14) and a subgroup of \(\widehat{G}\) is open if and only if it is closed and of finite index, or equivalently if its orthogonal is finite (Corollary 2 of II, p.~233); Corollary 4 of II, p.~228 shows that \((\widehat{\mathrm{G}})_{0}^{\perp}\) is the union of the finite subgroups of G, which is none other than \(B_{G}\) since G is discrete. We conclude by duality.

COROLLARY 1. --- Suppose G is compact. Then the following conditions are equivalent:

\begin{enumerate}
\def\labelenumi{(\roman{enumi})}
\item
  The group G is connected;
\item
  The group \(\widehat{\mathbf{G}}\) is torsion-free;
\item
  The group G is divisible.
\end{enumerate}

This follows from prop. 4 and cor. 7 of II, p.~229, since \(B_{G} = G\) .

COROLLARY 2. --- Suppose G is compact. Then G is totally disconnected if and only if \(\widehat{\mathbf{G}}\) is a torsion group.

The group G is totally disconnected if and only if its identity component is reduced to \(\{e\}\) ; prop. 4 shows that this condition is equivalent to \(B_{\widehat{G}} = \widehat{G}\) . Since the group \(\widehat{G}\) is discrete (prop. 18 of II, p.~233), this amounts to saying that each element of \(\widehat{G}\) generates a finite group, hence that \(\widehat{G}\) is torsion.

COROLLARY 3. --- If G is connected, then G is divisible.

Indeed, there exist a compact connected group K and an integer \(n \geqslant 0\) such that G is isomorphic to \(R^{n} \times K\) (II, p.~248, prop. 3, where one must have L = K). Corollary 1 shows that K is divisible, and therefore G is divisible.

\section{INVARIANT SUBSPACES}

The goal of this section is the study of certain subspaces invariant under translation in the spaces \(L^{1}(G)\) , \(L^{2}(G)\) and \(L^{\infty}(G)\) .

\subsection{\texorpdfstring{The case of Hilbert space \(L^{2}(G)\)}{The case of Hilbert space L\^{}\{2\}(G)}}

For every measurable subset M of \(\widehat{\mathbf{G}}\) , we denote by \(\mathrm{E_M}\) the set of \(f \in \mathrm{L}^2(\mathbf{G})\) such that the Fourier transform \(\mathcal{F}_{\mathrm{G}}(f)\) is zero almost everywhere on \(\widehat{\mathbf{G}}\) . Let \(\varphi_{\mathrm{M}}\) be the characteristic function of M. The space \(\mathrm{E_M}\) is the kernel of the continuous linear map \(f \mapsto \varphi_{\mathrm{M}}\mathcal{F}_{\mathrm{G}}(f)\) of \(\mathrm{L}^2(\mathbf{G})\) into \(\mathrm{L}^2(\widehat{\mathbf{G}})\) , and it is therefore a closed subspace of \(\mathrm{L}^2(\mathbf{G})\) .

PROPOSITION 1. --- a) Let M be a measurable subset of \(\widehat{\mathbf{G}}\) . For every \(x \in \mathbf{G}\) , the space \(\mathrm{E}_{\mathrm{M}}\) is stable under the map \(f \mapsto \varepsilon_x * f\) ;

\begin{enumerate}
\def\labelenumi{\alph{enumi})}
\setcounter{enumi}{1}
\item
  Let M and N be measurable subsets of G. One has \(E_{M} = E_{N}\) if and only if M and N are equal up to a locally negligible set;
\item
  Every subspace of \(L^{2}(G)\) stable under the maps \(f \mapsto \varepsilon_{x} * f\) for all \(x \in G\) is of the form \(E_{M}\) for some measurable subset M of \(\widehat{G}\) .
\end{enumerate}

This result will be proved later (cf.~V, to appear).
